\chapter{Polyhedral and tropical spaces}
\index{tropical space}\index{polyhedral space}
\label{chap:tropical-spaces}

Tropical spaces are assembled from weighted rational polyhedra. The weights
are not decoration: they are precisely what makes integration independent of
subdivision and makes boundary terms cancel. We therefore develop balancing
before introducing abstract charts.

\Cref{chap:superforms} produced a differential calculus on a real vector
space.  That calculus is affine invariant, but it lives on a vector space,
whereas the objects that occur in non-Archimedean geometry---skeleta of
models, tropicalizations of analytic domains, dual complexes---are neither
vector spaces nor manifolds.  They are locally finite unions of polyhedra, glued
along faces, and the gluing is exactly where the arithmetic enters: a cell
carries a lattice, a top-dimensional cell carries a multiplicity, and a
codimension-one wall carries a relation between the multiplicities of the
cells meeting along it.  The purpose of this chapter is to install that
combinatorial layer, and to prove that the superform calculus survives it.

\section{What the combinatorics must remember}

The underlying subset of a tropical cycle is not enough.  At a point of a
top-dimensional cell one must remember its tangent lattice and multiplicity;
at a codimension-one face one must remember primitive normal vectors.  These
data control, respectively, normalized volume, change of variables, and
Stokes cancellation.

The chapter has four goals:
\begin{enumerate}
\item pass cleanly between lattices, affine spaces, and polyhedral complexes;
\item formulate balancing intrinsically in a quotient lattice;
\item define smooth forms on a polyhedral support by restriction from an
  ambient affine space;
\item distinguish a weighted support, a calibration, and a tropical
  intersection multiplicity.
\end{enumerate}
The reader only needs elementary convex geometry and the superform calculus
of \Cref{chap:superforms}.  The stable-intersection argument is the
fan-displacement proof, written out in \Cref{thm:tropical-stable-intersection}.

\subsection*{What this chapter supplies to the rest of the book}

Five packages of results are assembled here and used repeatedly.

\begin{enumerate}
\item \emph{Polyhedral bookkeeping.}
\Cref{sec:ch10-polyhedra} fixes what a rational polyhedron, its tangent
lattice, its recession cone, its star, and its primitive normal vectors are,
and proves that two polyhedral structures on the same set admit a common
refinement (\Cref{prop:tropical-common-refinement}).  Every later
``after refining the two charts'' is an appeal to that proposition.

\item \emph{Balancing, and what it is equivalent to.}
\Cref{sec:ch10-balancing} defines weights and balancing and then proves that
balancing is equivalent to the vanishing of an explicit boundary map
(\Cref{thm:ch10-balancing-boundary}) and, in the one-dimensional
case, to the validity of Stokes' formula
(\Cref{thm:ch10-balancing-stokes-fan}); the role of balancing in
displacement-independence of the intersection product appears in
\Cref{ex:ch10-unbalanced} and in the proof of
\Cref{thm:tropical-stable-intersection}.  The Stokes equivalence is the
reason balancing appears at all: it is the exact hypothesis under which the
cancellation in \Cref{thm:balanced-stokes} takes place.

\item \emph{Forms on a support.}
\Cref{sec:ch10-forms} defines $A^{p,q}$ on an open subset of a polyhedral
support by restriction from an ambient affine space and proves that the
answer does not depend on that ambient space
(\Cref{thm:ch10-forms-presentation-independent}).  This is what allows
\Cref{chap:analytic-tropicalization} to define forms on a Berkovich space by
a limit over tropical charts.

\item \emph{Charts and calibrations.}
\Cref{sec:ch10-tropical-spaces,sec:ch10-calibrations} define tropical spaces
and calibrations and prove that an integer-affine transition map transports
the sheaf of superforms, the two differentials, the tangent lattices, and the
calibration (\Cref{thm:ch10-transition-transport}).  This is the structure
integrated in \Cref{chap:tropical-integration} and the structure carried by
the analytic tropicalizations of \Cref{chap:analytic-tropicalization}.

\item \emph{Stable intersection.}
\Cref{sec:ch10-stable-intersection,sec:ch10-gysin} construct the intersection
product of tropical cycles, first by hand in three examples, then in general
(\Cref{thm:tropical-stable-intersection}), and identify it with a Gysin
pullback along the diagonal
(\Cref{thm:ch10-diagonal-gysin}).  \Cref{chap:delta-forms} uses precisely
this product to multiply delta-forms
(\Cref{thm:delta-preform-algebra}).
\end{enumerate}

\subsection*{Dependencies}

\Cref{sec:ch10-polyhedra,sec:ch10-balancing} are pure convex geometry and
lattice algebra; they use nothing from earlier chapters except the polyhedral
vocabulary recorded in \Cref{sec:appG-polyhedra}.
\Cref{sec:ch10-forms} uses the superform algebra of
\Cref{chap:superforms}, in particular the pullback formula there and
\Cref{prop:superform-differential-identities}.
\Cref{sec:ch10-calibrations} uses the sign conventions of
\Cref{sec:appG-calibrations} and anticipates the integration of
\Cref{chap:tropical-integration}; a reader who wants to see the calibration
in action before reading its axioms may go directly to
\Cref{prop:integral-subdivision-independent}.
\Cref{sec:ch10-stable-intersection} is again self-contained convex geometry.
No analytic geometry is used anywhere in the chapter: the analytic origin of
the weights is the subject of \Cref{chap:analytic-tropicalization} and of
\Cref{thm:appG-analytic-calibration}.

The running toric example of \Cref{sec:ch01-running-toric} and the metric
graphs of \Cref{sec:metric-graphs} are the two smallest instances of
everything below, and they are used as consistency checks rather than
redeveloped.

\begin{warning}{Max-plus and valuation coordinates}
Mikhalkin--Rau use
$\mathbf T=\mathbf R\cup\{-\infty\}$ with tropical addition $\max$ and write
integer-affine functions in coordinates analogous to $\log|z_i|$.  Our
analytic chapters use valuation coordinates
$u_i=-\log|z_i|$.  Thus the translation is $x_i=-u_i$: max-plus polynomials
and min-plus valuation polynomials are related by
\[
                         F_{\max}(x)=-F_{\min}(-x).
\]
Thus changing the coordinate sign alone does \emph{not} exchange maximum
and minimum; one must also negate the value of the tropical polynomial.  We
use the Mikhalkin--Rau
max-plus terminology in this abstract tropical chapter and display the sign
change whenever an analytic tropicalization enters.  No geometric assertion
depends on choosing one of the two isomorphic semifields.
\end{warning}

The warning is not pedantry.  Both conventions occur in the sources this
chapter draws on, and the two errors they invite---forgetting the sign of the
coordinate, and forgetting the sign of the function---cancel each other in
some formulas and reinforce each other in others.  The safe procedure, used
throughout, is to record the tropical \emph{cycle} in whichever coordinate is
convenient, since a balanced complex remains balanced under $x\mapsto-x$
(the reflection is an integer-linear isomorphism), and to display the
translation explicitly whenever a tropical \emph{function} crosses between
the two pictures.

\section{Integer-affine polyhedra}
\label{sec:ch10-polyhedra}
\index{polyhedron!rational}\index{polyhedral complex}\index{primitive normal vector}\index{tangent lattice}

Let $N$ be a free abelian group of rank $r$, put
$N_{\mathbf R}=N\otimes_{\mathbf Z}\mathbf R$, and let
$M=\operatorname{Hom}(N,\mathbf Z)$.  We write
$\langle\cdot,\cdot\rangle$ for the pairing between $M$ and $N_{\mathbf R}$.
All polyhedra in this chapter are subsets of such an $N_{\mathbf R}$, or of
an affine space under it; the lattice is part of the data and is never
suppressed.

\begin{definition}
\label{def:ch10-integer-affine-polyhedra}
An \emph{integer-affine function} on $N_{\mathbf R}$ is $x\mapsto
\langle m,x\rangle+c$ with $m\in M$ and $c\in\mathbf R$. A
\emph{rational polyhedron} is an intersection of finitely many half-spaces
$\langle m_i,x\rangle\geq c_i$, with $m_i\in M$. A \emph{polyhedral complex}
$\mathscr C$ is a locally finite collection of polyhedra, closed under faces,
such that the intersection of two members is a face of each.
\end{definition}

This follows the terminology of Mikhalkin--Rau
\cite[\S\S1.3 and 2.1]{MR18Tropical}: 'rational' refers to the integral, or
equivalently rational, linear parts of the bounding inequalities; the
constants may be arbitrary real numbers.  CLD often says \emph{integral
affine} where tropical geometry says \emph{integer affine}; the two phrases
are synonyms in this book.  When every $c_i$ lies in a subgroup
$\Gamma\subset\mathbf R$, we add the qualifier $\Gamma$-rational.  None of
the balancing arguments below depends on the constants.

Following the same reference, a \emph{facet} of a pure polyhedral complex
means a top-dimensional cell.  For a face of codimension one in an individual
polyhedron we always write \emph{codimension-one face}; this prevents the two
common meanings of 'facet' from being confused.

A \emph{face} of a polyhedron $\sigma$ is a subset of the form
$\sigma\cap\{\langle m,\cdot\rangle=c\}$ for an integer covector $m$ with
$\langle m,\cdot\rangle\geq c$ on $\sigma$, together with $\sigma$ itself.
Faces of rational polyhedra are rational polyhedra, and the faces of a
polyhedron form a finite lattice under inclusion; we use these two facts
without further comment, referring to \cite[\S1.3]{MR18Tropical} or to any
standard account of convexity such as \cite{Roc70,Sch14}.  The
\emph{relative interior} $\RelInt(\sigma)$ is the interior of $\sigma$ inside
its affine span; a polyhedron is the disjoint union of the relative interiors
of its faces.

The \emph{recession cone} of a polyhedron $\sigma$ is
\[
 \operatorname{rec}(\sigma)
 =\{v\in N_{\mathbf R}\mid x+tv\in\sigma
           \text{ for every }x\in\sigma,\ t\geq0\}.
\]
The affine span of $\sigma$ is $x_0+N_{\sigma,\mathbf R}$, where
$N_\sigma=N\cap\operatorname{span}_{\mathbf R}(\sigma-\sigma)$ is saturated.
The quotient by $N_\sigma$ records transverse integral directions.

Two elementary facts about this lattice are used constantly, and we record
them once.

\begin{lemma}
\label{lem:ch10-tangent-lattice}
Let $\sigma\subset N_{\mathbf R}$ be a rational polyhedron of dimension $d$
and let $N_\sigma=N\cap\operatorname{span}_{\mathbf R}(\sigma-\sigma)$.
\begin{enumerate}
\item $N_\sigma$ is a saturated subgroup of $N$ of rank $d$, and
      $N_{\sigma,\mathbf R}=\operatorname{span}_{\mathbf R}(\sigma-\sigma)$.
\item $N/N_\sigma$ is a free abelian group of rank $r-d$, and
      $\bigwedge^dN_\sigma$ is free of rank one.
\item If $\tau$ is a face of $\sigma$, then $N_\tau\subseteq N_\sigma$ and
      $N_\sigma/N_\tau$ is saturated in $N/N_\tau$.
\end{enumerate}
\end{lemma}

\begin{proof}
(1)  By definition $N_\sigma$ is the intersection of $N$ with a real subspace,
hence saturated: if $kv\in N_\sigma$ with $k\geq1$ and $v\in N$, then $v$ lies
in the same real subspace and hence in $N_\sigma$.  Rationality of $\sigma$
means that its affine span is cut out by equations $\langle m_i,\cdot\rangle
=c_i$ with $m_i\in M$; therefore
$\operatorname{span}_{\mathbf R}(\sigma-\sigma)$ is the real span of the
rational subspace $\{v:\langle m_i,v\rangle=0\}$, and a rational subspace is
spanned by its lattice points.  Hence $N_\sigma$ has rank $d$ and spans it.

(2)  A quotient of a finitely generated free abelian group by a saturated
subgroup is torsion free, hence free.  The rank count is additive, and the
top exterior power of a free module of rank $d$ is free of rank one.

(3)  $N_\tau\subseteq N_\sigma$ because $\tau-\tau\subseteq\sigma-\sigma$.
If $k\bar v\in N_\sigma/N_\tau$ with $\bar v\in N/N_\tau$ and $k\geq1$, lift
$\bar v$ to $v\in N$; then $kv\in N_\sigma+N_\tau=N_\sigma$, so $v\in
N_\sigma$ by (1) and $\bar v\in N_\sigma/N_\tau$.
\end{proof}

If $\tau$ is a codimension-one face of $\sigma$, then by
\Cref{lem:ch10-tangent-lattice} the
quotient $N_\sigma/N_\tau$ is free of rank one. Its primitive generator pointing
from $\tau$ into $\sigma$ is denoted $u_{\sigma/\tau}$.  Explicitly, $u_{\sigma/\tau}$
is the unique generator of $N_\sigma/N_\tau$ represented by a vector
$v\in N_\sigma$ with $x+\varepsilon v\in\sigma$ for some (equivalently, for
every) $x\in\RelInt(\tau)$ and all small $\varepsilon>0$.  We refer to it as
the \emph{primitive normal} of $\sigma$ along $\tau$, or as the
\emph{inward} primitive vector; the corresponding outward vector is
$-u_{\sigma/\tau}$, and \Cref{chap:tropical-integration} uses the outward
convention for boundary traces.  The two conventions are compared in the
warning at \eqref{eq:appG-primitive-normal}.

\begin{example}[Primitive means lattice primitive]
\label{ex:ch10-primitive-ray}
In $N=\mathbf Z^2$, let $\tau$ be the origin and let $\sigma$ be the ray
generated by $(2,2)$.  Its tangent lattice is not
$\mathbf Z(2,2)$ but the saturation $\mathbf Z(1,1)$.  Hence
$u_{\sigma/\tau}=(1,1)$.  Using $(2,2)$ would double every boundary
contribution and would destroy invariance under a change of parametrization.
\end{example}

\begin{example}[A primitive normal is a quotient class]
\label{ex:ch10-primitive-quotient}
In $N=\mathbf Z^2$ let $\tau=\mathbf R(1,1)$ and let
$\sigma=\{(x,y):y\geq x\}$, so that $\tau$ is a codimension-one face of
$\sigma$.  Then $N_\tau=\mathbf Z(1,1)$ and
$N_\sigma=\mathbf Z^2$, so $N_\sigma/N_\tau\cong\mathbf Z$ and the class of
$(0,1)$ is a generator pointing into $\sigma$; hence
$u_{\sigma/\tau}=\overline{(0,1)}$.  The class of $(-1,0)$ is the same
element of the quotient, because $(0,1)-(-1,0)=(1,1)\in N_\tau$.  There is
therefore no preferred representative in $N$, and any formula that uses one
must be checked to be insensitive to the ambiguity.  The Stokes traces
\eqref{eq:superform-boundary-traces} are so checked: contracting with a
vector tangent to $\tau$ kills the form after restriction.
\end{example}

\begin{definition}
\label{def:ch10-star}
For a face $\tau$ of a complex $\mathscr C$, the \emph{star} of $\tau$ is the
fan in $N_{\mathbf R}/N_{\tau,\mathbf R}$ obtained by projecting the cones of
local directions of all $\sigma\supset\tau$.  It carries the weights of the
corresponding maximal cells.
\end{definition}

Concretely, fix $x\in\RelInt(\tau)$.  For $\sigma\supset\tau$ the set
\[
 \operatorname{cone}_x(\sigma)
 =\{v\in N_{\mathbf R}\mid x+\varepsilon v\in\sigma
    \text{ for all sufficiently small }\varepsilon>0\}
\]
is a rational polyhedral cone containing $N_{\tau,\mathbf R}$, and it does
not depend on the choice of $x\in\RelInt(\tau)$.  The image of
$\operatorname{cone}_x(\sigma)$ in $N_{\mathbf R}/N_{\tau,\mathbf R}$ is the
cone of $\operatorname{Star}_{\mathscr C}(\tau)$ attached to $\sigma$; its
tangent lattice is $N_\sigma/N_\tau$, which is saturated in $N/N_\tau$ by
\Cref{lem:ch10-tangent-lattice}.  The star is a fan because the local cones
of two cells meet in the local cone of their intersection.  Note that
$\dim\operatorname{Star}(\tau)=\dim\mathscr C-\dim\tau$: passing to the star
trades a cell for a cone through the origin and lowers the dimension by
exactly the dimension of $\tau$.

Balancing at a codimension-one cell says exactly that the one-dimensional weighted star has
zero vector sum.  Formulating the condition in the quotient makes it
independent of any choice of a complementary subspace.  This is worth
emphasizing, because the alternative---choosing a complement to
$N_{\tau,\mathbf R}$ and reading the normal vectors there---produces a
condition that appears to depend on the complement and, worse, tempts one to
normalize by a Euclidean length.  Neither the complement nor a Euclidean
structure exists canonically, and the balancing equation must not use one.

\begin{definition}
\label{def:ch10-subdivision}
A polyhedral complex $\mathscr C'$ \emph{subdivides} $\mathscr C$, and
$\mathscr C'$ is called a \emph{refinement} of $\mathscr C$, if
$|\mathscr C'|=|\mathscr C|$ and every cell of $\mathscr C'$ is contained in
a cell of $\mathscr C$.  Two complexes with the same support are
\emph{equivalent} if they admit a common refinement.
\end{definition}

Because a cell of $\mathscr C'$ has a relative interior, and because the
relative interiors of the cells of $\mathscr C$ partition $|\mathscr C|$,
each cell $\sigma'\in\mathscr C'$ is contained in a unique smallest cell of
$\mathscr C$, namely the one whose relative interior meets
$\RelInt(\sigma')$.  We call it the \emph{carrier} of $\sigma'$.  Two
subdivisions of the same complex need not be comparable, which is why
\Cref{prop:tropical-common-refinement} is needed.

\begin{lemma}
\label{lem:ch10-subdivision-dimension}
Let $\mathscr C'$ subdivide $\mathscr C$, with both complexes pure of
dimension $n$.  Every facet of $\mathscr C'$ has a facet of $\mathscr C$ as
its carrier.  For each facet $\sigma$ of $\mathscr C$, the relative interiors
of the facets $\sigma'$ with carrier $\sigma$ cover $\RelInt(\sigma)$ up to a
locally finite union of cells of dimension less than $n$.
\end{lemma}
\begin{proof}
Every cell of the subdivision lies in its carrier and the support is
unchanged.  Near any point of $\RelInt(\sigma)$ only finitely many cells of
$\mathscr C'$ occur.  An open subset of the $n$-dimensional affine space
$\operatorname{Aff}(\sigma)$ cannot be covered by finitely many proper
affine subspaces, so the $n$-cells of $\mathscr C'$ contained in $\sigma$
cover $\RelInt(\sigma)$ away from their proper faces.  Their carriers have
dimension $n$ and hence are facets of $\mathscr C$.  The same local argument
at every point gives the asserted locally finite exceptional union.
\end{proof}

\begin{proposition}
\label{prop:tropical-common-refinement}
Let $\mathscr C_1$ and $\mathscr C_2$ be finite polyhedral complexes in the
same affine space.  The faces of the nonempty intersections
$\sigma_1\cap\sigma_2$ generate a common polyhedral refinement on
$|\mathscr C_1|\cap|\mathscr C_2|$.  If both complexes have integer slopes,
so does the refinement.
\end{proposition}
\begin{proof}
An intersection of finitely many affine half-spaces is a polyhedron, and an
intersection of two such polyhedra retains integral normal directions.
Taking all faces makes the collection face-closed.  If two intersections
meet, their intersection is obtained by imposing simultaneously the active
equalities from both descriptions; after passage to the face-closed
collection it is a common face.  Finiteness is immediate.  The inclusion of
each new cell in an old cell gives the two refinement maps.
\end{proof}

For locally finite complexes the same intersection construction is locally
finite: a neighbourhood meeting finitely many cells of each complex meets
only finitely many pairs, and each intersection has finitely many faces.
This is the elementary combinatorial reason that
forms or integrals represented in two charts can be compared after
refinement; analytic tropical charts require the product construction of
\Cref{chap:analytic-tropicalization} in addition.

The next lemma is the reason that subdivision does not disturb any of the
lattice data.  It is used silently in almost every proof of the chapter, and
it is the precise content of the slogan ``subdivision changes cells, not
lattices''.

\begin{lemma}
\label{lem:ch10-subdivision-lattices}
Let $\mathscr C'$ subdivide $\mathscr C$ and let $\sigma'\in\mathscr C'$ have
carrier $\sigma\in\mathscr C$.  Then
$\dim\sigma'=\dim\sigma$ implies $N_{\sigma'}=N_\sigma$.  If moreover $\tau'$
is a codimension-one face of $\sigma'$ whose relative interior meets
$\RelInt(\tau)$ for a codimension-one face $\tau$ of $\sigma$, then
$N_{\tau'}=N_\tau$ and $u_{\sigma'/\tau'}=u_{\sigma/\tau}$ under the
resulting identification $N_{\sigma'}/N_{\tau'}=N_\sigma/N_\tau$.
\end{lemma}

\begin{proof}
If $\dim\sigma'=\dim\sigma$ and $\sigma'\subseteq\sigma$, the two affine
spans coincide, since a $d$-dimensional convex set spans a $d$-dimensional
affine space and $\RelInt(\sigma')$ is open in that span.  Hence
$N_{\sigma'}=N\cap\operatorname{span}_{\mathbf R}(\sigma'-\sigma')
=N\cap\operatorname{span}_{\mathbf R}(\sigma-\sigma)=N_\sigma$.  The same
argument applied to $\tau'\subseteq\tau$ gives $N_{\tau'}=N_\tau$, because
$\RelInt(\tau')$ meets $\RelInt(\tau)$ and both have dimension
$\dim\sigma-1$.  Finally $u_{\sigma'/\tau'}$ is characterized by the property
that it is the generator of $N_{\sigma'}/N_{\tau'}$ represented by a vector
pointing from $\RelInt(\tau')$ into $\sigma'$.  Near a point of
$\RelInt(\tau')\cap\RelInt(\tau)$ the sets $\sigma'$ and $\sigma$ have the
same germ on the $\tau$ side, so the two characterizations select the same
generator.
\end{proof}

\begin{example}[Two incomparable subdivisions]
\label{ex:ch10-incomparable}
Let $\mathscr C$ consist of the unit square $[0,1]^2\subset\mathbf R^2$ and
its faces.  Cutting along the diagonal $y=x$ gives a subdivision
$\mathscr C_1$ into two triangles; cutting along $y=1-x$ gives another
subdivision $\mathscr C_2$.  Neither refines the other, and their common
refinement, produced by \Cref{prop:tropical-common-refinement}, consists of
the four small triangles obtained by cutting along both diagonals.  All of
these complexes have the same support and the same tangent lattice
$\mathbf Z^2$ on every two-dimensional cell, illustrating
\Cref{lem:ch10-subdivision-lattices}.  The new interior edges are the ones
whose contributions will have to cancel in
\Cref{thm:balanced-stokes}.
\end{example}

\section{Weights and balancing}
\label{sec:ch10-balancing}
\index{balancing condition}\index{tropical cycle}\index{weight!on a polyhedral cell}

A weighted complex is a polyhedral complex together with one integer per
top-dimensional cell.  Balancing is a linear condition on those integers,
imposed at every codimension-one cell.  We give the definition, then three
proofs that it is the right condition: it is the vanishing of a boundary
(\Cref{thm:ch10-balancing-boundary}); it is exactly what makes Stokes'
formula hold (\Cref{thm:ch10-balancing-stokes-fan}); and it is the hypothesis
making the stable intersection product of
\Cref{sec:ch10-stable-intersection} independent of the displacement
direction. The unbalanced example shows how independence can fail.

\begin{definition}
\label{def:ch10-weighted-balanced}
A pure $n$-dimensional weighted complex $(\mathscr C,m)$ assigns an integer
$m_\sigma\neq0$ to every facet $\sigma$.  We also write
$\omega(\sigma)=m_\sigma$, as in \cite[\S4.2]{MR18Tropical}. It is
\emph{balanced} if for every
$(n-1)$-cell $\tau$,
\[
                  \sum_{\sigma\supset\tau}
                    m_\sigma u_{\sigma/\tau}=0
                    \quad\text{in }N/N_\tau.
\]
Two weighted complexes are \emph{equivalent} if they have a common weighted
refinement.  A \emph{tropical $n$-cycle} is an equivalence class of balanced
pure $n$-dimensional weighted rational polyhedral complexes.  A tropical
cycle represented by positive weights is a \emph{tropical subspace}; a cycle
represented by a fan is a \emph{tropical fan cycle}.
\end{definition}

Thus negative coefficients are allowed in the group of tropical cycles.
The analytic tropicalizations used later carry positive multiplicities and
hence define tropical subspaces.  Cells of weight zero are omitted from the
support.  These conventions agree with
\cite[Definitions~4.2.1--4.2.3]{MR18Tropical}.

Four remarks on the definition are worth making before any example.

First, the sum in the balancing equation is finite, because a locally finite
complex has only finitely many facets containing a given cell.

Second, the equation is imposed on \emph{every} codimension-one cell, including
those that lie in the topological boundary of the support.  A complex with a
free edge---say the single segment $[0,1]\subset\mathbf R$ with weight
one---is therefore \emph{not} balanced: at the vertex $0$ there is one
incoming facet and the sum is the nonzero class $u_{[0,1]/\{0\}}=+1$.  This
is not a defect of the definition.  Balanced complexes are the objects
without boundary, and the boundary theory is developed separately in
\Cref{chap:tropical-integration}.

Third, weights are attached to facets, not to points.  Two facets may share a
point without sharing a weight, and the value of the weight function at a
point of a codimension-one cell is not defined.

Fourth, the support of a weighted complex, as a topological space, does not
determine the weights.  A tropical cycle is the pair, and two cycles with the
same support and different weights are different cycles.  The following
statement records the additive structure that this remark makes possible.

\begin{proposition}
\label{prop:ch10-cycle-group}
Fix $n$ and $N$.  The set $Z_n(N_{\mathbf R})$ of tropical $n$-cycles in
$N_{\mathbf R}$ represented by finite complexes, together with the zero
cycle, is an abelian group under the
following operation: given finite representatives $(\mathscr C_1,m_1)$ and
$(\mathscr C_2,m_2)$, choose a polyhedral overlay $\mathscr C$ on
$|\mathscr C_1|\cup|\mathscr C_2|$ whose restriction to each support
refines the corresponding complex, and put
\[
 (m_1+m_2)_\sigma
 =\begin{cases}
   m_{1,\sigma}+m_{2,\sigma},&\RelInt(\sigma)\subseteq|\mathscr C_1|\cap|\mathscr C_2|,\\
   m_{1,\sigma},&\RelInt(\sigma)\subseteq|\mathscr C_1|\setminus|\mathscr C_2|,\\
   m_{2,\sigma},&\RelInt(\sigma)\subseteq|\mathscr C_2|\setminus|\mathscr C_1|,
  \end{cases}
\]
discarding the cells of weight zero.  The inverse of $(\mathscr C,m)$ is
$(\mathscr C,-m)$.
\end{proposition}

\begin{proof}
\emph{Step 1: a complex covering both supports.}
Choose finite representatives $\mathscr C_1$ and $\mathscr C_2$.  Each of
their cells is an intersection of finitely many rational half-spaces, and
each affine hull of a cell is an intersection of finitely many rational
hyperplanes.  Let $\mathscr H$ be the finite set of all rational hyperplanes
occurring in such descriptions, and let $\mathscr A$ be the polyhedral
complex of $N_{\mathbf R}$ whose cells are the closures of the nonempty sets
obtained by prescribing, for each $H\in\mathscr H$, one of the three
conditions ``on $H$'', ``strictly on one side'', ``strictly on the other
side''.  This is a rational polyhedral complex: the sets in question are
relatively open convex polyhedra, their closures are rational polyhedra
whose faces are again of the same form, and two of them meet in a common
face.  By construction every cell of $\mathscr C_1$ and of $\mathscr C_2$ is
a union of cells of $\mathscr A$.  Let $\mathscr C$ be the subcomplex of
$\mathscr A$ consisting of the cells contained in
$|\mathscr C_1|\cup|\mathscr C_2|$, together with their faces, and discard
cells not contained in an $n$-cell; then $\mathscr C$ is pure of dimension
$n$, refines $\mathscr C_1$ over $|\mathscr C_1|$ and $\mathscr C_2$ over
$|\mathscr C_2|$, and has support $|\mathscr C_1|\cup|\mathscr C_2|$.

\emph{Step 2: the weights are well defined and balanced.}
A facet $\sigma$ of $\mathscr C$ has $\RelInt(\sigma)$ contained in
$|\mathscr C_1|$ or in its complement, and likewise for $|\mathscr C_2|$;
hence exactly one of the three cases of the displayed formula occurs, and
$m_{i,\sigma}$ means the weight of the unique facet of $\mathscr C_i$
containing $\sigma$, which exists by
\Cref{lem:ch10-subdivision-dimension}.  Extend $m_i$ by zero to all facets of
$\mathscr C$, obtaining $\widehat m_i$.  We claim $\widehat m_i$ is balanced
on $\mathscr C$.  Let $\tau$ be a codimension-one cell of $\mathscr C$.  If
$\tau\not\subseteq|\mathscr C_i|$, then no facet of $\mathscr C$ containing
$\tau$ lies in $|\mathscr C_i|$, since such a facet would contain $\tau$;
so the balancing sum is empty.  If $\tau\subseteq|\mathscr C_i|$, the facets
of $\mathscr C$ containing $\tau$ and contained in $|\mathscr C_i|$ are
exactly the facets of the refinement of $\mathscr C_i$ containing $\tau$, and
the balancing sum is the one furnished by
\Cref{prop:ch10-balancing-subdivision}.  Thus
$\widehat m_1+\widehat m_2$ is balanced, by
\Cref{cor:ch10-balancing-consequences}\,\cref{itm:ch10-bal-linear}, and it is
the weighting displayed in the statement.

\emph{Step 3: group axioms.}
Associativity, commutativity, and the inverse assertion are immediate on a
common refinement of all complexes involved, produced by Step 1 applied
repeatedly.  Discarding weight-zero cells changes neither the balancing
equations, in which those cells contribute nothing, nor the equivalence
class; and the result is independent of the refinement used, because two
refinements admit a common further refinement on which the two answers agree
facet by facet.
\end{proof}

\begin{example}[The one-dimensional vertex equation]
\label{ex:ch10-vertex-equation}
Let $u_1,\ldots,u_s$ be the primitive outgoing directions of a weighted
one-dimensional complex at a vertex.  Since the tangent lattice of the
vertex is zero,
balancing is simply
\[
                         \sum_{j=1}^s m_j u_j=0
                         \quad\text{in }N.
\]
For two rays with positive weights this forces them to lie on the same line
with opposite primitive directions and equal weights.  A genuine branch
point therefore needs at least three directions.
\end{example}

\begin{example}[The tropical line]
\label{ex:ch10-tropical-line}
The three rays in $\mathbf R^2$ generated by
$(1,0)$, $(0,1)$, and $(-1,-1)$, each with weight one, are balanced at the
origin because their primitive vectors sum to zero. Rays contained in an
open half-plane cannot form a balanced one-cycle with positive weights:
pairing with the defining covector $m$ gives
$\langle m,\sum_km_ku_k\rangle>0$.  Rays in a \emph{closed} half-plane
can: $\mathbf R_{\geq0}(1,0)\cup\mathbf R_{\geq0}(-1,0)$ with unit
weights is the line $\mathbf Re_1$, subdivided at the origin.
\end{example}

\begin{figure}[ht]
\centering
\begin{tikzpicture}[x=1cm,y=1cm,>=Stealth]
  \coordinate (xzero) at (-3,0);
  \coordinate (uzero) at (3,0);

  \node[font=\small] at (-3,2.62) {max-plus coordinates $x$};
  \draw[-{Stealth[length=2.2mm]},very thick] (xzero)--(-0.85,0)
    node[pos=0.82,above,font=\small] {$(1,0)$};
  \draw[-{Stealth[length=2.2mm]},very thick] (xzero)--(-3,2.15)
    node[above,font=\small] {$(0,1)$};
  \draw[-{Stealth[length=2.2mm]},very thick] (xzero)--(-4.55,-1.55)
    node[below left,font=\small] {$(-1,-1)$};
  \fill (xzero) circle (1.5pt);
  \node[above right,font=\small] at (-2.9,0.08) {$0$};

  \node[font=\small] at (3,2.62) {min-plus valuation coordinates $u$};
  \draw[-{Stealth[length=2.2mm]},very thick] (uzero)--(0.85,0)
    node[pos=0.82,below,font=\small] {$(-1,0)$};
  \draw[-{Stealth[length=2.2mm]},very thick] (uzero)--(3,-2.15)
    node[below,font=\small] {$(0,-1)$};
  \draw[-{Stealth[length=2.2mm]},very thick] (uzero)--(4.55,1.55)
    node[above right,font=\small] {$(1,1)$};
  \fill (uzero) circle (1.5pt);
  \node[above right,font=\small] at (3.08,0.08) {$0$};

  \draw[{Stealth[length=2mm]}-{Stealth[length=2mm]},thin]
    (-0.82,2.05)--(0.82,2.05)
    node[midway,above,font=\small] {$x=-u$};
  \node[font=\scriptsize,fill=white,inner sep=1pt] at (-3,-2.32)
    {$F_{\max}(x)=\max\{0,-x_1,-x_2\}$};
  \node[font=\scriptsize,fill=white,inner sep=1pt] at (3,-2.72)
    {$F_{\min}(u)=\min\{0,-u_1,-u_2\}$};
  \node[font=\scriptsize] at (0,-1.45)
    {$F_{\max}(x)=-F_{\min}(-x)$};
\end{tikzpicture}
\caption{The tropical line and the convention firewall.  On the left the
three weight-one primitive directions sum to zero.  The change $x=-u$ sends
them to the three directions on the right, whose sum is again zero; negating
the function as well gives $F_{\max}(x)=-F_{\min}(-x)$.}
\label{fig:tropical-line-balancing}
\end{figure}

With the max convention, the corner locus of
$F_{\max}(x)=\max\{0,-x_1,-x_2\}$ is the left-hand cycle in
\Cref{fig:tropical-line-balancing}.  In valuation coordinates the paired
min-plus function is
$F_{\min}(u)=\min\{0,-u_1,-u_2\}$; both the coordinate and the function value
are negated, exactly as prescribed in the warning above.

\begin{proposition}
\label{prop:ch10-balancing-subdivision}
Let $\mathscr C$ be pure dimensional.  Balancing is preserved by subdivision
when each new facet inherits the weight of its carrier.
\end{proposition}
\begin{proof}
Let $\mathscr C'$ subdivide $\mathscr C$, of pure dimension $n$, and let
$\tau'$ be a codimension-one cell of $\mathscr C'$.  Its carrier in
$\mathscr C$ has dimension either $n$ or $n-1$.

If the carrier of $\tau'$ is an $n$-cell $\sigma$ of $\mathscr C$, then a
neighbourhood of a point of $\RelInt(\tau')$ inside $|\mathscr C|$ is an
$n$-dimensional convex set contained in $\sigma$, cut by $\tau'$ into two
pieces.  Exactly two facets $\sigma_1',\sigma_2'$ of $\mathscr C'$ contain
$\tau'$, both with carrier $\sigma$, hence both with weight $m_\sigma$ and
tangent lattice $N_\sigma$; and
$u_{\sigma_1'/\tau'}=-u_{\sigma_2'/\tau'}$ because the two germs lie on
opposite sides of the hyperplane $\operatorname{span}(\tau'-\tau')$ in
$N_{\sigma,\mathbf R}$.  The balancing sum at $\tau'$ is therefore
$m_\sigma(u_{\sigma_1'/\tau'}+u_{\sigma_2'/\tau'})=0$.

If instead the carrier of $\tau'$ is a codimension-one cell $\tau$ of
$\mathscr C$, then $N_{\tau'}=N_\tau$ by
\Cref{lem:ch10-subdivision-lattices}, and the facets $\sigma'\supset\tau'$
are partitioned according to their carriers $\sigma\supset\tau$; for each
$\sigma$ there is exactly one such $\sigma'$ near a fixed point of
$\RelInt(\tau')$, with $m_{\sigma'}=m_\sigma$ and
$u_{\sigma'/\tau'}=u_{\sigma/\tau}$.  The two balancing sums coincide term by
term.
\end{proof}

\subsection{Balancing as the vanishing of a boundary}
\label{sec:ch10-balancing-boundary}

The balancing equation was written above as a relation among primitive normal
vectors.  That is the form one computes with, but it hides the reason the
condition is stable under everything in sight.  We now show that balancing is
the vanishing of an explicit incidence boundary attached to $\mathscr C$.
It is the top-degree part of the cellular boundary with coefficients in
$\bigwedge^nN$; the coefficient of a facet $\sigma$ is required to lie in
the line $\bigwedge^nN_\sigma$. We use locally finite chains so that an
infinite locally finite complex may carry nonzero weights on every facet.

Fix a pure $n$-dimensional rational polyhedral complex $\mathscr C$ in
$N_{\mathbf R}$.  For $k\leq n$ let $\mathscr C_k$ denote its set of
$k$-cells.  Choose, once and for all, an orientation of every cell of
dimension $n$ and $n-1$, that is, a generator
\[
                  e_\sigma\in\textstyle\bigwedge^{\dim\sigma}N_\sigma,
\]
which is possible by \Cref{lem:ch10-tangent-lattice}\,(2).  For a facet
$\sigma$ and a codimension-one face $\tau\prec\sigma$ define the
\emph{incidence number} $[\sigma:\tau]\in\{\pm1\}$ by
\begin{equation}
\label{eq:ch10-incidence}
 \widetilde u_{\sigma/\tau}\wedge e_\tau
 =-[\sigma:\tau]\,e_\sigma
 \quad\text{in }\textstyle\bigwedge^nN_\sigma,
\end{equation}
where $\widetilde u_{\sigma/\tau}\in N_\sigma$ is any lift of the primitive
normal.  The left-hand side does not depend on the lift, since two lifts
differ by an element $w\in N_\tau$ and $w\wedge e_\tau=0$, the element
$e_\tau$ being a top exterior power of $N_\tau$; and it is a generator of
$\bigwedge^nN_\sigma$,
because a lift $\widetilde u_{\sigma/\tau}$ together with a lattice basis of
$N_\tau$ representing $e_\tau$ is a lattice basis of $N_\sigma$.  The sign in
\eqref{eq:ch10-incidence} is chosen so that $[\sigma:\tau]$ is the usual
incidence number for the \emph{outward} normal $-u_{\sigma/\tau}$, matching
the boundary convention of \Cref{chap:tropical-integration}.

\begin{definition}
\label{def:ch10-calibrated-chain}
Set
\[
 C_n(\mathscr C)=\prod_{\sigma\in\mathscr C_n}\textstyle\bigwedge^nN_\sigma,
 \qquad
 C_{n-1}(\mathscr C)=\prod_{\tau\in\mathscr C_{n-1}}\textstyle\bigwedge^nN,
\]
and define $\partial:C_n(\mathscr C)\to C_{n-1}(\mathscr C)$ by
\begin{equation}
\label{eq:ch10-boundary}
 (\partial\mu)_\tau
 =\sum_{\sigma\supset\tau}[\sigma:\tau]\,\mu_\sigma
 \qquad(\tau\in\mathscr C_{n-1}),
\end{equation}
using the inclusions $\bigwedge^nN_\sigma\subset\bigwedge^nN$.  Given weights
$m$, the associated \emph{calibrated chain} is
$\mu(m)=(m_\sigma e_\sigma)_{\sigma\in\mathscr C_n}$.
\end{definition}

Every component of the boundary is a finite sum by local finiteness.
Thus the map is defined on these products; for a finite complex they are
also finite direct sums. No finite-support condition is imposed on $m$.

The inclusions used in \eqref{eq:ch10-boundary} are genuinely injective:
$N_\sigma$ is saturated in $N$ by \Cref{lem:ch10-tangent-lattice}, so
$N/N_\sigma$ is free and $N\cong N_\sigma\oplus(N/N_\sigma)$; taking exterior
powers exhibits $\bigwedge^nN_\sigma$ as a direct summand of
$\bigwedge^nN$.  This is where saturation is used, and it is the reason the
convention of \Cref{sec:appG-polyhedra}---compute lattice data only after
saturating---cannot be relaxed.

\begin{theorem}
\label{thm:ch10-balancing-boundary}
Let $(\mathscr C,m)$ be a pure $n$-dimensional weighted rational polyhedral
complex in $N_{\mathbf R}$.  Then
\[
 (\mathscr C,m)\text{ is balanced}
 \quad\Longleftrightarrow\quad
 \partial\mu(m)=0 .
\]
More precisely, for every codimension-one cell $\tau$ one has
\begin{equation}
\label{eq:ch10-boundary-vs-balancing}
 (\partial\mu(m))_\tau
 =-\Bigl(\sum_{\sigma\supset\tau}m_\sigma\widetilde u_{\sigma/\tau}\Bigr)
   \wedge e_\tau,
\end{equation}
and the map $N/N_\tau\to\bigwedge^nN$, $\bar v\mapsto\widetilde v\wedge
e_\tau$, is injective.  The conclusion does not depend on the chosen
orientations $e_\sigma,e_\tau$.
\end{theorem}

\begin{proof}
\emph{Step 1: the displayed formula.}
Fix $\tau\in\mathscr C_{n-1}$ and a facet $\sigma\supset\tau$.  By
\eqref{eq:ch10-incidence},
\[
 [\sigma:\tau]\,e_\sigma
 =-\widetilde u_{\sigma/\tau}\wedge e_\tau .
\]
Multiplying by $m_\sigma$ and summing over $\sigma\supset\tau$ gives
\eqref{eq:ch10-boundary-vs-balancing}.  Note that the right-hand side is
computed in $\bigwedge^nN$, where all the individual terms live, and that it
does not depend on the choice of lifts $\widetilde u_{\sigma/\tau}$, since
changing one of them alters the bracket by an element of $N_\tau$, which
dies against $e_\tau$.

\emph{Step 2: injectivity of $\bar v\mapsto\widetilde v\wedge e_\tau$.}
The map $N\to\bigwedge^nN$, $v\mapsto v\wedge e_\tau$, is $\mathbf Z$-linear
and kills $N_\tau$, hence descends to $N/N_\tau$.  Suppose
$\widetilde v\wedge e_\tau=0$ with $\widetilde v\in N$.  Extend $e_\tau$ to
$\bigwedge^{n-1}N_{\tau,\mathbf R}$ and work in $\bigwedge^nN_{\mathbf R}$,
where the same relation holds.  Choosing a real basis $f_1,\ldots,f_{n-1}$ of
$N_{\tau,\mathbf R}$ with $e_\tau=f_1\wedge\cdots\wedge f_{n-1}$, the
vanishing of $\widetilde v\wedge f_1\wedge\cdots\wedge f_{n-1}$ in the
exterior algebra of $N_{\mathbf R}$ forces
$\widetilde v\in\operatorname{span}_{\mathbf R}(f_1,\ldots,f_{n-1})
=N_{\tau,\mathbf R}$.  Since $N_\tau$ is saturated
(\Cref{lem:ch10-tangent-lattice}), $\widetilde v\in N_\tau$, that is,
$\bar v=0$.

\emph{Step 3: conclusion.}
By Steps 1 and 2, $(\partial\mu(m))_\tau=0$ if and only if
$\sum_{\sigma\supset\tau}m_\sigma u_{\sigma/\tau}=0$ in $N/N_\tau$, which is
the balancing equation at $\tau$.  Ranging over $\tau$ gives the
equivalence.

\emph{Step 4: independence of orientations.}
Replacing $e_\sigma$ by $-e_\sigma$ replaces $\mu(m)_\sigma$ by its negative
and, by \eqref{eq:ch10-incidence}, $[\sigma:\tau]$ by its negative for every
$\tau$; the product $[\sigma:\tau]\mu(m)_\sigma$ is unchanged.  Replacing
$e_\tau$ by $-e_\tau$ replaces every $[\sigma:\tau]$ with $\sigma\supset\tau$
by its negative, hence multiplies $(\partial\mu(m))_\tau$ by $-1$ and does
not affect its vanishing.
\end{proof}

The formulation \eqref{eq:ch10-boundary} is not merely cosmetic.  Three
consequences that were proved by hand above, or will be needed below, are
now visible at once.

\begin{corollary}
\label{cor:ch10-balancing-consequences}
Let $(\mathscr C,m)$ be a pure $n$-dimensional weighted rational polyhedral
complex.
\begin{enumerate}
\item\label{itm:ch10-bal-linear} The balanced weightings form a subgroup of
      $\prod_{\sigma\in\mathscr C_n}\mathbf Z$; in particular sums and
      integer multiples of balanced weightings are balanced.
\item\label{itm:ch10-bal-local} $(\mathscr C,m)$ is balanced if and only if
      $\operatorname{Star}_{\mathscr C}(\tau)$, with the inherited weights,
      is balanced for every cell $\tau$ of $\mathscr C$ of dimension at most
      $n-1$.
\item\label{itm:ch10-bal-affine} If $\mathscr C$ consists of a single
      $n$-dimensional affine subspace and its faces, then every weighting is
      balanced, because $\mathscr C_{n-1}=\varnothing$.
\end{enumerate}
\end{corollary}

\begin{proof}
\cref{itm:ch10-bal-linear} The chain $\mu(m)$ depends
$\mathbf Z$-linearly on $m$ once the orientations are fixed, and $\partial$
is $\mathbf Z$-linear.

\cref{itm:ch10-bal-local} It suffices to treat $\dim\tau=n-1$, since a
codimension-one cell of $\operatorname{Star}(\tau)$ for $\dim\tau<n-1$
corresponds to a codimension-one cell $\tau'$ of $\mathscr C$ containing
$\tau$, and the two balancing equations agree.  For $\dim\tau=n-1$ the star
is the fan in $N_{\mathbf R}/N_{\tau,\mathbf R}$ whose rays are the classes
$\mathbf R_{\geq0}u_{\sigma/\tau}$, with tangent lattices
$N_\sigma/N_\tau\subset N/N_\tau$; the balancing equation of the star at its
unique zero-dimensional cell is literally the balancing equation of
$\mathscr C$ at $\tau$.

\cref{itm:ch10-bal-affine} An affine subspace, viewed as a complex with
one $n$-cell and no proper faces of dimension $n-1$ inside the complex, has
an empty set of codimension-one cells, so the condition is vacuous.
\end{proof}

The reader should note that
\Cref{cor:ch10-balancing-consequences}\,\cref{itm:ch10-bal-affine} is a
statement about the \emph{complex}, not about the set: if one artificially
subdivides an affine subspace, codimension-one cells appear, and
\Cref{prop:ch10-balancing-subdivision} guarantees that the equations at those
new cells hold automatically.  This is the first place where it matters that
balancing is a property of the equivalence class and not of a chosen
polyhedral structure.

\begin{warning}{The boundary complex is not a homology theory here}
We use only the single map \eqref{eq:ch10-boundary}.  Extending the picture
to a full chain complex, with $\partial^2=0$ and coefficient systems varying
by dimension, is the starting point of tropical homology.  Nothing below
requires that theory, and we do not develop it.  What we do use is precisely
that the map is well defined after orientations are chosen, and that its
vanishing is independent of those choices by Step~4 above.
\end{warning}

\subsection{Balancing as Stokes' formula}
\label{sec:ch10-balancing-stokes}

The chain-level statement explains why balancing is stable.  The next
theorem explains why one wants it at all: on a one-dimensional fan,
balancing is \emph{equivalent} to the vanishing of the integral of an exact
top-degree form.  This is the whole mechanism of
\Cref{thm:balanced-stokes}, isolated in the smallest case where it can be
computed in closed form.

\begin{theorem}
\label{thm:ch10-balancing-stokes-fan}
Let $u_1,\ldots,u_s\in N$ be primitive vectors spanning pairwise distinct
rays, let $m_1,\ldots,m_s$ be nonzero integers, and let
$X=\bigcup_k\mathbf R_{\geq0}u_k$ with the weights $m_k$.  For a superform
$\eta=\sum_{i}a_i\,\ddprime x_i\in A^{0,1}_c(N_{\mathbf R})$ with compact
support, write $a(0)=(a_1(0),\ldots,a_r(0))$.  Then
\begin{equation}
\label{eq:ch10-fan-stokes}
 \int_X\dprime\eta
 =-\Bigl\langle a(0),\ \sum_{k=1}^sm_ku_k\Bigr\rangle,
\end{equation}
where the left-hand side is the calibrated integral of
\Cref{chap:tropical-integration}.  Consequently $X$ is balanced if and only
if $\int_X\dprime\eta=0$ for every compactly supported
$\eta\in A^{0,1}$.
\end{theorem}

\begin{proof}
Fix $k$ and parametrize the ray $\mathbf R_{\geq0}u_k$ by $s\mapsto su_k$.
This is an integer-affine isomorphism from $\mathbf R_{\geq0}$, with its
standard lattice $\mathbf Z$, onto the ray with its tangent lattice
$\mathbf Z u_k$, because $u_k$ is primitive.  Hence $s$ is a dual affine
coordinate attached to an oriented lattice basis of that tangent lattice, and
the definition of the calibrated integral in
\Cref{chap:tropical-integration} reads it as $m_k$ times the ordinary
integral in the variable $s$; that this does not depend on the chosen basis
is \Cref{prop:integral-subdivision-independent}.

Pulling back, the primed and double-primed one-forms become
$\dprime x_i\mapsto u_{k,i}\dprime s$ and
$\ddprime x_i\mapsto u_{k,i}\ddprime s$.  Hence
\[
 \dprime\eta
 =\sum_{i,j}\frac{\partial a_i}{\partial x_j}\,
   \dprime x_j\wedge\ddprime x_i
\]
pulls back to
\[
 \Bigl(\sum_{i,j}\frac{\partial a_i}{\partial x_j}(su_k)\,
   u_{k,j}u_{k,i}\Bigr)\dprime s\wedge\ddprime s
 =\frac{\rd}{\rd s}\Bigl[\sum_ia_i(su_k)u_{k,i}\Bigr]\,
   \dprime s\wedge\ddprime s,
\]
by the chain rule.  The bracket is $\langle a(su_k),u_k\rangle$.  Since
$\eta$ has compact support, the fundamental theorem of calculus gives
\[
 \int_{\mathbf R_{\geq0}u_k}\dprime\eta
 =\int_0^{\infty}\frac{\rd}{\rd s}
   \langle a(su_k),u_k\rangle\,\rd s
 =-\langle a(0),u_k\rangle .
\]
Multiplying by $m_k$ and summing over $k$ gives
\eqref{eq:ch10-fan-stokes}.

If $X$ is balanced then $\sum_km_ku_k=0$, since the tangent lattice of the
vertex is zero, and the right-hand side of \eqref{eq:ch10-fan-stokes}
vanishes for every $\eta$.  Conversely, if the right-hand side vanishes for
every $\eta$, choose $\eta=\chi\,\ddprime x_i$ with $\chi$ compactly
supported and equal to $1$ near the origin; then $a(0)$ is the $i$th standard
covector and \eqref{eq:ch10-fan-stokes} gives the $i$th coordinate of
$\sum_km_ku_k$.  Letting $i$ vary gives $\sum_km_ku_k=0$.
\end{proof}

The identity \eqref{eq:ch10-fan-stokes} identifies the vector
$\sum_km_ku_k$ precisely as the obstruction to Stokes' formula.
Chambert--Loir and Ducros build this obstruction into the axioms of their
calibrated theory under the name of harmony; the identification of harmony
with balancing in the analytic setting is recorded in
\Cref{thm:appG-analytic-calibration}.

\subsection{A gallery of fans}
\label{sec:ch10-fan-gallery}

We now compute.  The examples in this subsection are the ones referred to
throughout the rest of \Cref{part:cld}, and each of them is verified
completely.  Throughout, $e_1,\ldots,e_r$ is the standard basis of
$N=\mathbf Z^r$ and we set
\[
                 e_0=-(e_1+\cdots+e_r),
\]
so that $e_0+e_1+\cdots+e_r=0$.

\begin{example}[The standard tropical line in $\mathbf R^r$]
\label{ex:ch10-standard-line}
Let $L_r\subset\mathbf R^r$ be the union of the $r+1$ rays
$\mathbf R_{\geq0}e_i$, $i=0,\ldots,r$, each with weight one.  It is a
one-dimensional fan, and it is balanced: the tangent lattice of its unique
vertex is zero, and $\sum_{i=0}^re_i=0$ by construction.  For $r=2$ this is
the picture of \Cref{fig:tropical-line-balancing}; for $r=1$ it is the whole
line $\mathbf R$, subdivided at the origin.

For $r=2$, and only for $r=2$, the fan $L_r$ is the corner locus, in the
sense of \Cref{def:appG-corner-locus}, of the max-plus polynomial
\[
 F_{\max}(x)=\max\{0,-x_1,\ldots,-x_r\}.
\]
Indeed, the maximum in $F_{\max}$ is attained twice on a set of pure
codimension one, whereas $\dim L_r=1$; the two dimensions agree exactly when
$r=2$.  For $r=2$ the corner locus is the union of
\[
 \{x_1=0\leq x_2\},\qquad
 \{x_2=0\leq x_1\},\qquad
 \{x_1=x_2\leq0\},
\]
which are the three rays $\mathbf R_{\geq0}e_2$,
$\mathbf R_{\geq0}e_1$, and $\mathbf R_{\geq0}e_0$.  For $r\geq3$ the corner
locus is the standard tropical hyperplane, of which
\Cref{ex:ch10-standard-plane} is the case $r=3$, and $L_r$ has codimension
$r-1$.

Deleting any one of the $r+1$ rays destroys balancing: the remaining $r$
primitive vectors are linearly independent, so no nonzero integer
combination of them vanishes.  Deleting a ray and reweighting the others
does not help either, for the same reason.
\end{example}

\begin{example}[The standard tropical plane in $\mathbf R^3$]
\label{ex:ch10-standard-plane}
Let
\[
 H=\bigcup_{0\leq i<j\leq3}\operatorname{cone}(e_i,e_j)
 \subset\mathbf R^3,
\]
with $e_0=-(e_1+e_2+e_3)$ and weight one on each of the six
two-dimensional cones.  Then $H$ is a balanced two-dimensional fan, and it is
the corner locus of $F_{\max}(x)=\max\{0,-x_1,-x_2,-x_3\}$.

\emph{Balancing.}  The codimension-one cells of $H$ are the four rays
$\rho_i=\mathbf R_{\geq0}e_i$.  The facets containing $\rho_i$ are the three
cones $\operatorname{cone}(e_i,e_j)$ with $j\neq i$.  For $i\neq j$ the pair
$\{e_i,e_j\}$ extends to a lattice basis of $\mathbf Z^3$: for
$i,j\geq1$ this is clear, and for $i=0$, $j=1$ one computes
\[
 \bigl|\det\bigl(e_0,e_1,e_2\bigr)\bigr|
 =\bigl|\det\bigl((-1,-1,-1),(1,0,0),(0,1,0)\bigr)\bigr|=1,
\]
the remaining cases following by permuting the coordinates.  Hence the
tangent lattice of $\operatorname{cone}(e_i,e_j)$ is
$\mathbf Ze_i+\mathbf Ze_j$, this group being saturated, and
the primitive normal along $\rho_i$ is the class of $e_j$ in
$\mathbf Z^3/\mathbf Ze_i$.  The balancing equation at $\rho_i$ therefore
reads
\[
 \sum_{j\neq i}\overline{e_j}
 =\overline{\textstyle\sum_{j\neq i}e_j}
 =\overline{-e_i}=0
 \quad\text{in }\mathbf Z^3/\mathbf Ze_i,
\]
which holds.  Note how the computation used only the relation
$\sum_ie_i=0$; the same argument proves balancing of the analogous fan in
$\mathbf R^r$ in every codimension.

\emph{Corner locus.}  On the cone $\operatorname{cone}(e_2,e_3)$ one has
$x_1=0$ and $x_2,x_3\geq0$, so $0=-x_1\geq-x_2,-x_3$ and the maximum is
attained twice; on $\operatorname{cone}(e_0,e_3)$ one has $x_1=x_2\leq0$ and
$x_3\geq x_1$, so $-x_1=-x_2\geq0,-x_3$.  The remaining four cones are
obtained by permuting indices.  Conversely, a point at which the maximum is
attained twice satisfies one of these two systems.
\end{example}

The next example is the first one whose balancing is not visible from a
tropical polynomial, because the fan has codimension two.

\begin{example}[A balanced fan of codimension two]
\label{ex:ch10-codim-two-fan}
In $N=\mathbf Z^3$ consider the five primitive vectors
\[
\begin{aligned}
 w_1&=(1,0,0), & w_2&=(0,1,0), & w_3&=(0,0,1),\\
 w_4&=(1,-1,0), & w_5&=(-2,0,-1),
\end{aligned}
\]
and let $W=\bigcup_{k=1}^5\mathbf R_{\geq0}w_k$ with all weights equal to
one.  Each $w_k$ is primitive: the greatest common divisor of its
coordinates is one.  The five rays are pairwise distinct, and
\[
 w_1+w_2+w_3+w_4+w_5
 =(1+0+0+1-2,\,0+1+0-1+0,\,0+0+1+0-1)
 =(0,0,0),
\]
so $W$ is balanced at its unique vertex.

$W$ is not the corner locus of a tropical polynomial, for the trivial reason
that such a locus has codimension one in $\mathbf R^3$ whereas $\dim W=1$.
It is also indecomposable, in the following precise sense: $W$ is not the
sum, in the group $Z_1(\mathbf R^3)$ of
\Cref{prop:ch10-cycle-group}, of two nonzero effective cycles supported on
subsets of its rays.  Indeed, in such a decomposition the two weightings are
positive integers adding up to $1$ on each ray, so each ray belongs to
exactly one summand, with weight one; the summands are then the fans on a
partition of $\{w_1,\ldots,w_5\}$ into two nonempty subfamilies, and
balancing of each summand says that each subfamily sums to zero.  A direct
check eliminates this.  No two of the $w_k$ are
negatives of each other, and no $w_k$ is zero.  For triples one computes
\[
\begin{aligned}
 w_1{+}w_2{+}w_3&=(1,1,1), &
 w_1{+}w_2{+}w_4&=(2,0,0), &
 w_1{+}w_2{+}w_5&=(-1,1,-1),\\
 w_1{+}w_3{+}w_4&=(2,-1,1), &
 w_1{+}w_3{+}w_5&=(-1,0,0), &
 w_1{+}w_4{+}w_5&=(0,-1,-1),\\
 w_2{+}w_3{+}w_4&=(1,0,1), &
 w_2{+}w_3{+}w_5&=(-2,1,0), &
 w_2{+}w_4{+}w_5&=(-1,0,-1),\\
 w_3{+}w_4{+}w_5&=(-1,-1,0), & & & &
\end{aligned}
\]
none of which vanishes, and each four-element subfamily sums to the negative
of the remaining $w_k$, hence is nonzero.

So $W$ is a genuinely five-valent balanced curve in $\mathbf R^3$, not
assembled from lines and not read off a tropical polynomial.  Whether $W$ is
the tropicalization of an algebraic curve in the three-dimensional torus over
some valued field is a separate and much harder question, to which balancing
gives no answer; \Cref{chap:analytic-tropicalization} proves that
tropicalizations are balanced, not that balanced cycles are
tropicalizations.
\end{example}

\begin{example}[A weighted fan that fails balancing, and what goes wrong]
\label{ex:ch10-unbalanced}
In $N=\mathbf Z^2$ let $X=\mathbf R_{\geq0}e_1\cup\mathbf R_{\geq0}e_2$, with
weight one on each ray.  The balancing defect at the origin is
$e_1+e_2=(1,1)\neq0$, so $X$ is not balanced.  Three separate failures
follow, and it is instructive to see all three, because they are the three
places where balancing is used later.

\emph{Stokes fails.}  Take $\eta=\chi\,\ddprime x_1$ with $\chi$ smooth,
compactly supported, and equal to $1$ near the origin.  By
\eqref{eq:ch10-fan-stokes},
\[
 \int_X\dprime\eta
 =-\langle(1,0),\,e_1+e_2\rangle=-1\neq0 .
\]
Thus the current of integration over $X$ is not $\dprime$-closed. A
boundary-free Green formula of the type of
\Cref{cor:polyhedral-integration-by-parts} needs an additional vertex term
on $X$.

\emph{The calibrated chain is not a cycle.}  With
$e_{\mathbf R_{\geq0}e_1}=e_1$, $e_{\mathbf R_{\geq0}e_2}=e_2$ and
$e_{\{0\}}=1\in\bigwedge^0N$, the boundary
\eqref{eq:ch10-boundary} at the vertex is $-(e_1+e_2)\neq0$, in accordance
with \Cref{thm:ch10-balancing-boundary}.

\emph{Stable intersection is not well defined.}  Let $Y=\mathbf Re_2$, the
whole vertical line, with weight one; it is balanced.  Displace $Y$ by
$\varepsilon v$ with $v=(v_1,v_2)$ and $\varepsilon>0$ small.  The displaced
line is $\{x_1=\varepsilon v_1\}$.  If $v_1>0$, it meets $X$ in the single
point $(\varepsilon v_1,0)$, lying in the relative interior of
$\mathbf R_{\geq0}e_1$; the intersection is transverse and the lattice index
$[\mathbf Z^2:\mathbf Ze_1+\mathbf Ze_2]$ equals one, so the count is $1$.
If $v_1<0$, the displaced line meets neither ray and the count is $0$.  The
two answers disagree, so no displacement-independent product exists.
Comparing with the wall-crossing step in the proof of
\Cref{thm:tropical-stable-intersection}: there $H=N_{\tau,\mathbf R}+
N'_{\mathbf R}=\mathbf Re_2$ is a hyperplane, the classes of $e_1$ and $e_2$
in $N/(N\cap H)\cong\mathbf Z$ are $\lambda_1=1$ and $\lambda_2=0$, and the
quantity $\sum_k m_k\lambda_k$, which balancing would force to vanish, equals
$1$.  The discrepancy between the two counts is exactly that number.
\end{example}

\begin{example}[Balancing with unequal weights]
\label{ex:ch10-weights-234}
In $\mathbf Z^2$ take the rays generated by $(1,0)$, $(0,1)$, $(-2,-3)$ with
weights $2$, $3$, $1$.  Then
\[
 2(1,0)+3(0,1)+1\cdot(-2,-3)=(0,0),
\]
so the fan is balanced.  All three vectors are primitive, and no reweighting
by a common factor is possible, so this is not a multiple of the tropical
line.  The example shows that the weight of a facet is not determined by its
tangent lattice: the ray $\mathbf R_{\geq0}(1,0)$ carries weight $2$ here and
weight $1$ in \Cref{ex:ch10-standard-line}, with the same lattice
$\mathbf Z(1,0)$.
\end{example}

\subsection{Bergman fans of matroids}
\label{sec:ch10-bergman}
\index{Bergman fan}\index{matroid}

The last family of examples is the standard source of balanced fans that are
not visibly tropicalizations of hypersurfaces, and it is the one the reader
is most likely to meet again.  We record the definition, verify balancing in
the case used as an anchor, and are explicit about what we do not prove.

Let $E=\{1,\ldots,n\}$ and let $M$ be a loopless matroid on $E$ of rank $\rho$, given
by its rank function $\operatorname{rk}$.  A \emph{flat} of $M$ is a subset
$F\subseteq E$ with $\operatorname{rk}(F\cup\{i\})>\operatorname{rk}(F)$ for
every $i\notin F$.  In $\mathbf R^E$ write $e_S=\sum_{i\in S}e_i$ and work in
the quotient
\[
 N_M=\mathbf Z^E/\mathbf Ze_E,
 \qquad
 N_{M,\mathbf R}=\mathbf R^E/\mathbf Re_E,
\]
a lattice of rank $n-1$; the quotient is by a primitive vector, so $N_M$ is
free.

\begin{definition}
\label{def:ch10-bergman-fan}
The \emph{Bergman fan} $B(M)\subset N_{M,\mathbf R}$ is the fan whose cones
are
\[
 \sigma_{\mathscr F}
 =\operatorname{cone}
   \bigl(\overline{e_{F_1}},\ldots,\overline{e_{F_k}}\bigr),
\]
one for each chain $\mathscr F:\varnothing\subsetneq F_1\subsetneq\cdots
\subsetneq F_k\subsetneq E$ of proper nonempty flats of $M$, together with
the zero cone.  Every maximal cone is given weight one.
\end{definition}

Since a chain of proper nonempty flats has at most $\rho-1$ members---the
ranks $\operatorname{rk}(F_1)<\cdots<\operatorname{rk}(F_k)$ are distinct
integers strictly between $0$ and $\rho$---the fan $B(M)$ has dimension at
most $\rho-1$, and it is pure of dimension exactly $\rho-1$ when the maximal
chains are the complete flags of flats.  We use only the rank-two case,
where everything can be checked by hand.

\begin{proposition}
\label{prop:ch10-bergman-rank-two}
Let $M$ be a matroid of rank two on $E$ with no loops, and let
$F_1,\ldots,F_s$ be its rank-one flats.  Then
$B(M)=\bigcup_{k}\mathbf R_{\geq0}\overline{e_{F_k}}$ is a one-dimensional
fan in $N_{M,\mathbf R}$, the vectors $\overline{e_{F_k}}$ are primitive,
and $B(M)$ is balanced with all weights equal to one.
\end{proposition}

\begin{proof}
A chain of proper nonempty flats of a rank-two matroid consists of a single
rank-one flat, so the cones of $B(M)$ are the rays
$\mathbf R_{\geq0}\overline{e_{F_k}}$ together with $\{0\}$; the fan is
one-dimensional, that is, of dimension $\rho-1=1$.

The rank-one flats of a loopless matroid partition $E$: every $i\in E$ lies
in the closure of $\{i\}$, which is a rank-one flat since $i$ is not a loop,
and two rank-one flats that meet are equal because a common element already
has rank one and flats of equal rank containing a common rank-one set
coincide.  Hence
\[
 \sum_{k=1}^se_{F_k}=e_E,
 \qquad\text{so}\qquad
 \sum_{k=1}^s\overline{e_{F_k}}=0
 \quad\text{in }N_M .
\]
Since the tangent lattice of the vertex $\{0\}$ is zero, this is precisely
the balancing equation.

It remains to check primitivity of $\overline{e_{F}}$ for a rank-one flat
$F$.  Suppose $\overline{e_F}=d\,\overline v$ with $v\in\mathbf Z^E$ and
$d\geq1$; then $e_F=dv+ce_E$ for some $c\in\mathbf Z$.  Reading coordinate
$i\in F$ gives $1=dv_i+c$ and coordinate $j\notin F$ gives $0=dv_j+c$; such
$j$ exists because $F$ is a proper subset, $M$ having rank two.  Subtracting,
$1=d(v_i-v_j)$, so $d=1$.
\end{proof}

\begin{example}[The Bergman fan of $U_{2,3}$ is the tripod]
\label{ex:ch10-bergman-u23}
Let $M=U_{2,3}$ be the uniform matroid of rank two on $E=\{1,2,3\}$: every
subset of size at most two is independent, and the bases are the three
two-element subsets.  Its flats are $\varnothing$, the three singletons, and
$E$; so the proper nonempty flats are $\{1\},\{2\},\{3\}$ and
\Cref{prop:ch10-bergman-rank-two} applies with $s=3$.  Thus
\[
 B(U_{2,3})
 =\mathbf R_{\geq0}\overline{e_1}\cup
  \mathbf R_{\geq0}\overline{e_2}\cup
  \mathbf R_{\geq0}\overline{e_3}
 \subset\mathbf Z^3/\mathbf Ze_E\otimes\mathbf R,
\]
a one-dimensional balanced fan with three rays: the tripod.  Choosing the
lattice isomorphism $\mathbf Z^3/\mathbf Ze_E\cong\mathbf Z^2$ that sends
$\overline{e_1}\mapsto(1,0)$ and $\overline{e_2}\mapsto(0,1)$, one gets
$\overline{e_3}=-\overline{e_1}-\overline{e_2}\mapsto(-1,-1)$, so
$B(U_{2,3})$ is carried to the standard tropical line $L_2$ of
\Cref{ex:ch10-standard-line}.

Two sanity checks.  The dimension is $\rho-1=1$, as it must be for a rank-two
matroid.  And $U_{2,3}$ is the matroid of three pairwise independent vectors
spanning a plane, for instance the columns of
$\left(\begin{smallmatrix}1&0&-1\\0&1&-1\end{smallmatrix}\right)$; the
corresponding projective linear space is the line
$\{z_1+z_2+z_3=0\}$ in $\mathbf P^2$, intersected with its dense torus.
Over a trivially valued field its tropicalization in valuation coordinates
is $B(U_{2,3})$: every circuit minimum is attained at least twice. Reflection
to max-plus coordinates uses the paired sign convention above. So in this
case the Bergman fan is visibly a tropicalization.
\end{example}

\begin{example}[A Bergman fan that is a subdivided plane]
\label{ex:ch10-bergman-boolean}
Let $M$ be the Boolean matroid of rank three on $E=\{1,2,3\}$, that is, the
matroid of a basis of a three-dimensional space.  Every subset is a flat, so
the proper nonempty flats are the three singletons and the three pairs, and
the maximal chains are the six flags $\{i\}\subsetneq\{i,j\}$.  Hence $B(M)$
has six two-dimensional cones inside the rank-two lattice
$\mathbf Z^3/\mathbf Ze_E$.  Under the isomorphism of
\Cref{ex:ch10-bergman-u23}, the six rays are
\[
 \overline{e_1}=(1,0),\quad
 \overline{e_2}=(0,1),\quad
 \overline{e_3}=(-1,-1),
\]
\[
 \overline{e_{\{1,2\}}}=(1,1),\quad
 \overline{e_{\{2,3\}}}=(-1,0),\quad
 \overline{e_{\{1,3\}}}=(0,-1),
\]
and the six cones are
\[
\begin{gathered}
 \operatorname{cone}((1,0),(1,1)),\quad
 \operatorname{cone}((0,1),(1,1)),\quad
 \operatorname{cone}((0,1),(-1,0)),\\
 \operatorname{cone}((-1,-1),(-1,0)),\quad
 \operatorname{cone}((-1,-1),(0,-1)),\quad
 \operatorname{cone}((1,0),(0,-1)),
\end{gathered}
\]
listed by the angular sectors they occupy.
Their union is all of $\mathbf R^2$, and consecutive cones share exactly one
ray.  Thus $B(M)$ is a subdivision of the plane with weight one, which is
balanced by \Cref{prop:ch10-balancing-subdivision} and
\Cref{cor:ch10-balancing-consequences}\,\cref{itm:ch10-bal-affine}.  The
example is a useful check on the dimension formula: a rank-three matroid
gives a two-dimensional fan.
\end{example}

The general fact behind these computations is that $B(M)$ is a balanced pure
$(\rho-1)$-dimensional fan for every loopless matroid $M$ of rank $\rho$.
When $M$ is realizable over a field, its Bergman fan is the tropicalization
of the corresponding linear space.  These results are due to Ardila--Klivans
and Sturmfels; they are not proved here and are not used below.  The precise
conclusion needed for orientation is only that balancing is weaker than
realizability as a tropical \emph{linear} space: a nonrealizable matroid still
has a balanced Bergman fan, but it cannot arise from a linear realization of
that matroid.

\begin{warning}{Subdivision versus pushforward}
Under subdivision, a top-dimensional cell keeps its tangent lattice and its
weight.  Lattice-index factors enter a different operation: pushforward under
an integer-affine map.  Mixing the two rules is a common source of incorrect
multiplicities.
\end{warning}

\section{Forms on a polyhedral support}
\label{sec:ch10-forms}
\index{superform!on a polyhedral support}

A polyhedral support need not be a manifold. Its topology alone does not
specify a smooth sheaf. We use the calculus of \Cref{chap:superforms} by
restricting forms from the ambient affine space. The resulting algebra is
unchanged under the integer-affine presentation changes stated in
the theorem below. That proof is
\Cref{thm:ch10-forms-presentation-independent}, and it is the technical
heart of the passage from tropical charts to Berkovich spaces in
\Cref{chap:analytic-tropicalization}.

Let $|\mathscr C|$ be the support. A smooth superform on an open subset
$\Omega\subset|\mathscr C|$ is represented by a smooth superform on a
neighborhood in $N_{\mathbf R}$.

\begin{definition}
\label{def:ch10-forms-on-support}
Let $W\subset N_{\mathbf R}$ be open with
$W\cap|\mathscr C|=\Omega$, and set
\begin{equation}
\label{eq:ch10-form-kernel}
\begin{gathered}
 K^{p,q}(W,\mathscr C;\Omega)
 =\bigl\{\alpha\in A^{p,q}(W)\ \bigm|\
   \alpha_x(v_1,\ldots,v_p;w_1,\ldots,w_q)=0\bigr\},\\
 \text{imposed for all }\sigma\in\mathscr C,\
   x\in\Omega\cap\sigma
   \text{ and } v_i,w_j\in N_{\sigma,\mathbf R},
\end{gathered}
\end{equation}
where the final parameter will be omitted when $\Omega$ is clear.  Define
\begin{equation}
\label{eq:ch10-forms-direct-limit}
 A^{p,q}(\Omega)
 =\varinjlim_{\substack{W\subset N_{\mathbf R}\text{ open}\\
                         W\cap|\mathscr C|=\Omega}}
   A^{p,q}(W)/K^{p,q}(W,\mathscr C;\Omega),
\end{equation}
the transition maps being restrictions to smaller ambient neighbourhoods.
Thus a form is a germ of an ambient form along $\Omega$.  Two representatives
are equal if, after restriction to a common ambient neighbourhood, their
difference vanishes on tangent vectors to every cell. Equivalently, at each
point of a cell its value lies in the exterior-algebra ideal
generated by covectors annihilating that cell's tangent space.
\end{definition}

Two points deserve emphasis.  First, the condition is imposed cell by cell:
at a point
$x$ lying in several cells, one tests separately with tangent vectors of each
cell containing $x$, not with vectors from different cells simultaneously.
Second, the quotient depends only on the collection of pairs
$(x,N_{\sigma,\mathbf R})$ with $x\in\Omega\cap\sigma$, that is, on the
support together with its local tangent data.  In particular a subdivision of
$\mathscr C$ does not change $A^{p,q}(\Omega)$, since by
\Cref{lem:ch10-subdivision-lattices} the top-dimensional tangent spaces are
unchanged and the tangent space of a smaller cell is contained in that of a
cell of the original complex which contains it.

The quotient is local in $\Omega$.  A form may be represented on different
open subsets by different ambient extensions, provided their restrictions
agree on every cell. For the relative interior of a single $d$-dimensional
cell, coordinates along its span identify the quotient with the ordinary
superform algebra on an open subset of $\mathbf R^d$; normal covectors
disappear. At its boundary, smooth forms retain extension germs along the cell.

\begin{lemma}
\label{lem:ch10-forms-sheaf}
The assignment $\Omega\mapsto A^{p,q}(\Omega)$, with the evident restriction
maps, is a sheaf on $|\mathscr C|$.
\end{lemma}

\begin{proof}
Restriction maps exist because the kernel in
\eqref{eq:ch10-form-kernel} is defined pointwise and because restriction is
compatible with the direct limit in \eqref{eq:ch10-forms-direct-limit}.
Separatedness is the same pointwise observation: if a form vanishes on a
cover, one representative has zero tangential evaluation near every point of
$\Omega$, hence lies in the kernel along all of $\Omega$.

For gluing, let $\Omega=\bigcup_i\Omega_i$ and let $\alpha_i$ be forms on
$\Omega_i$ agreeing on overlaps.  Choose ambient representatives locally.
Since the polyhedral support is metrizable and paracompact, refine the cover
near $\Omega$ to a locally finite family of ambient open sets $(V_a)$ such
that $V_a\cap|\mathscr C|$ lies in some $\Omega_{i(a)}$ and
$\alpha_{i(a)}$ has an ambient representative
$\widetilde\alpha_a$ on $V_a$.  After shrinking once more, a smooth
partition of unity $(\chi_a)$ is subordinate to $(V_a)$ on an ambient
neighbourhood $V$ of $\Omega$, with $\operatorname{supp}(\chi_a)\subset
V_a$.  Extend $\chi_a\widetilde\alpha_a$ by zero to $V$ and set
$\widetilde\alpha=\sum_a\chi_a\widetilde\alpha_a$.  The sum is locally
finite.  At $x\in\Omega\cap\sigma$, all representatives occurring in the
sum have the same evaluation on tangent vectors of $\sigma$, because the
original sections agree on overlaps.  Since $\sum_a\chi_a(x)=1$,
$\widetilde\alpha$ has that common evaluation.  It therefore represents a
form whose restriction to every $\Omega_i$ is $\alpha_i$.
\end{proof}

The partition-of-unity argument is worth pausing on: it uses smooth bump
functions on the \emph{ambient} affine space, not on $|\mathscr C|$.  This is
the recurring pattern of the chapter.  Ambient smoothness is a tool;
the statements proved with it are about the support.

\subsection{Independence of the ambient presentation}
\label{sec:ch10-presentation}

A polyhedral support arises in practice as the image of a tropicalization
map, and different maps give different ambient spaces.  We must know that the
form algebra does not depend on the choice.  The precise statement is the
following.  Recall that an integer-affine map $F:N_{\mathbf R}\to
N'_{\mathbf R}$ has linear part carrying $N$ into $N'$ and an arbitrary real
translation, as in \Cref{def:appG-affine-conventions}.

\begin{theorem}
\label{thm:ch10-forms-presentation-independent}
Let $\mathscr C$ and $\mathscr C'$ be rational polyhedral complexes in
$N_{\mathbf R}$ and $N'_{\mathbf R}$, let $\Omega\subseteq|\mathscr C|$ and
$\Omega'\subseteq|\mathscr C'|$ be open, and let
$\varphi:\Omega\to\Omega'$ be a homeomorphism.  Assume:
\begin{enumerate}
\item\label{itm:ch10-pres-local} every point of $\Omega$ has a neighbourhood
      on which $\varphi$ is the restriction of an integer-affine map
      $N_{\mathbf R}\to N'_{\mathbf R}$, and every point of $\Omega'$ has a
      neighbourhood on which $\varphi^{-1}$ is the restriction of an
      integer-affine map $N'_{\mathbf R}\to N_{\mathbf R}$;
\item\label{itm:ch10-pres-cells} after refining $\mathscr C$ and
      $\mathscr C'$, the map $\varphi$ carries $\Omega\cap\sigma$ onto
      $\Omega'\cap\sigma'$ for a bijection $\sigma\mapsto\sigma'$ of the
      cells meeting $\Omega$ with those meeting $\Omega'$.
\end{enumerate}
Then $\varphi^*:A^{p,q}(\Omega')\to A^{p,q}(\Omega)$ is a well-defined
isomorphism of bigraded algebras, compatible with $\dprime$, $\ddprime$, and
the involution $J$.  Moreover the linear part of $\varphi$ induces
isomorphisms $N_{\sigma}\cong N'_{\sigma'}$ of tangent lattices.
\end{theorem}

\begin{proof}
\emph{Step 1: reduction to a local statement.}
Both sides are sheaves on $\Omega$ by \Cref{lem:ch10-forms-sheaf}, and the
construction of $\varphi^*$ below is local; so by
\Cref{lem:ch10-forms-sheaf} it suffices to construct compatible
isomorphisms over a cover of $\Omega$ by open sets on which $\varphi$ and
$\varphi^{-1}$ are restrictions of integer-affine maps.  Fix such an open
set $\Omega_0$ and integer-affine maps $F:N_{\mathbf R}\to N'_{\mathbf R}$
and $G:N'_{\mathbf R}\to N_{\mathbf R}$ with $F|_{\Omega_0}=\varphi$ and
$G|_{\varphi(\Omega_0)}=\varphi^{-1}$.

\emph{Step 2: $\varphi^*$ is well defined.}
The map $F$ is affine, so $F^*$ carries $A^{p,q}$ of a neighbourhood of
$\varphi(\Omega_0)$ in $N'_{\mathbf R}$ to $A^{p,q}$ of a neighbourhood of
$\Omega_0$ in $N_{\mathbf R}$, and commutes with wedge, $\dprime$,
$\ddprime$, and $J$; this is the pullback of \Cref{chap:superforms}.  We
must check that $F^*$ carries $K^{p,q}$ into $K^{p,q}$.  Let
$\alpha'\in K^{p,q}$ and let $x\in\Omega_0\cap\sigma$ with tangent vectors
$v_i,w_j\in N_{\sigma,\mathbf R}$.  Writing $DF$ for the linear part,
\[
 (F^*\alpha')_x(v_1,\ldots;w_1,\ldots)
 =\alpha'_{F(x)}(DF(v_1),\ldots;DF(w_1),\ldots).
\]
By hypothesis \cref{itm:ch10-pres-cells}, $F$ maps a neighbourhood of $x$ in
$\sigma$ into $\sigma'$.  The cone of directions $v$ with $x+tv\in\sigma$ for
small $t>0$ has nonempty interior in $N_{\sigma,\mathbf R}$, because $\sigma$
is convex of full dimension in its affine span and $x\in\sigma$; hence that
cone spans $N_{\sigma,\mathbf R}$.  For $v$ in the cone,
$DF(v)=\lim_{t\downarrow0}t^{-1}(F(x+tv)-F(x))$ lies in
$N'_{\sigma',\mathbf R}$, and by linearity
$DF(N_{\sigma,\mathbf R})\subseteq N'_{\sigma',\mathbf R}$.  Hence the
right-hand side vanishes, and $F^*\alpha'\in K^{p,q}$.

\emph{Step 3: $\varphi^*$ is independent of the chosen extension.}
If $F_1,F_2$ both restrict to $\varphi$ on $\Omega_0$, then for
$x\in\Omega_0\cap\sigma$ and $v\in N_{\sigma,\mathbf R}$ one has
$DF_1(v)=DF_2(v)$: this holds for $v$ in the spanning cone of feasible
directions of Step 2, because $F_1=F_2$ on $\Omega_0$ and both differ from
their common value at $x$ by the same difference quotients, and it then holds
on $N_{\sigma,\mathbf R}$ by linearity.  Thus $F_1^*\alpha'$ and
$F_2^*\alpha'$ have the same evaluations on cell tangent vectors at points of
$\Omega_0$, that is, they differ by an element of $K^{p,q}$.  Consequently
the induced map on quotients depends only on $\varphi$, and the local maps
agree on overlaps and glue.

\emph{Step 4: $\varphi^*$ is bijective.}
By Steps 2 and 3 applied to $\varphi^{-1}$ we get a map
$(\varphi^{-1})^*:A^{p,q}(\Omega)\to A^{p,q}(\Omega')$, and locally it is
induced by $G^*$.  Now $G\circ F$ restricts to the identity on $\Omega_0$,
so by Step 3, $(G\circ F)^*=F^*\circ G^*$ induces the identity on
$A^{p,q}(\Omega_0)$; symmetrically for $F\circ G$.  Hence $\varphi^*$ and
$(\varphi^{-1})^*$ are mutually inverse.

\emph{Step 5: tangent lattices.}
Fix a cell $\sigma$ meeting $\Omega$ and a point
$x\in\RelInt(\sigma)\cap\Omega$.  As in Step 2, $DF$ maps
$N_{\sigma,\mathbf R}$ into $N'_{\sigma',\mathbf R}$, and $DG$ maps back;
since $G\circ F$ and $F\circ G$ restrict to identities near $x$ and near
$F(x)$, the two linear maps are mutually inverse isomorphisms
$N_{\sigma,\mathbf R}\cong N'_{\sigma',\mathbf R}$.  Because $DF(N)\subseteq
N'$ and $DG(N')\subseteq N$, they restrict to mutually inverse maps
$N\cap N_{\sigma,\mathbf R}\to N'\cap N'_{\sigma',\mathbf R}$, that is, to an
isomorphism $N_\sigma\cong N'_{\sigma'}$.
\end{proof}

The theorem is stated for a homeomorphism because that is the form needed for
charts.  The case that occurs when one enlarges the ambient space---replacing
a tropicalization by a finer one with more coordinates---is worth extracting
separately, since it is the one used to compare two charts through their
product.

\begin{corollary}
\label{cor:ch10-forms-graph}
Let $\mathscr C_1,\mathscr C_2$ be rational polyhedral complexes in
$N_{1,\mathbf R}$, $N_{2,\mathbf R}$, let
$\Omega_i\subseteq|\mathscr C_i|$ be open, and let
$\varphi:\Omega_1\to\Omega_2$ satisfy the hypotheses of
\Cref{thm:ch10-forms-presentation-independent}.  Let
\[
 \Gamma=\{(x,\varphi(x))\mid x\in\Omega_1\}
 \subset N_{1,\mathbf R}\oplus N_{2,\mathbf R},
\]
with the polyhedral structure induced by
\cref{itm:ch10-pres-cells}, and let $p_1,p_2$ be the two projections.  Then
$p_1$ and $p_2$ restrict to homeomorphisms $\Gamma\to\Omega_1$ and
$\Gamma\to\Omega_2$ satisfying the same hypotheses, and
\[
 p_1^*:A^{p,q}(\Omega_1)\longrightarrow A^{p,q}(\Gamma),
 \qquad
 p_2^*:A^{p,q}(\Omega_2)\longrightarrow A^{p,q}(\Gamma)
\]
are isomorphisms with $p_1^*=p_2^*\circ\varphi_*$, where
$\varphi_*=(\varphi^{-1})^*$.
\end{corollary}

\begin{proof}
The projections are restrictions of the linear maps
$N_{1,\mathbf R}\oplus N_{2,\mathbf R}\to N_{i,\mathbf R}$, which are integer
linear; their inverses on $\Gamma$ are $x\mapsto(x,\varphi(x))$ and
$y\mapsto(\varphi^{-1}(y),y)$, which are locally integer affine because
$\varphi$ and $\varphi^{-1}$ are.  The cell correspondence is
$\sigma\leftrightarrow(\sigma\times\sigma')\cap\Gamma$.  Hence
\Cref{thm:ch10-forms-presentation-independent} applies to $p_1|_\Gamma$ and
$p_2|_\Gamma$ and gives the two isomorphisms.  The identity
$p_1|_\Gamma=\varphi^{-1}\circ p_2|_\Gamma$ gives the last assertion after
applying the functoriality established in Step 4 of that proof.
\end{proof}

The graph construction compares two presentations when their images are
related by the homeomorphism in the corollary. In general two
tropicalizations of one analytic domain need not be homeomorphic, and the
product moment map need not have a graph as its image.
\Cref{chap:analytic-tropicalization} uses product moment maps as common
refinements, with the additional analytic comparison results stated there.

\begin{proposition}
\label{prop:ch10-forms-quotient-operations}
Wedge products, $\dprime$, and $\ddprime$ descend to the quotient.  So does
pullback by an integer-affine map that carries each source cell into a target
cell.
\end{proposition}
\begin{proof}
If a covector annihilates a cell tangent space, its wedge with any form still
vanishes on tangent vectors, so wedge products descend.

For the differential assertion, let
$\alpha\in K^{p,q}(W,\mathscr C;\Omega)$; we must show
$\dprime\alpha\in K^{p+1,q}(W,\mathscr C;\Omega)$.  Fix a cell $\sigma$ and
first take $x\in\RelInt(\sigma)\cap\Omega$.  Choose affine coordinates
$y_1,\ldots,y_d,z_1,\ldots,z_{r-d}$ on $N_{\mathbf R}$ adapted to $\sigma$,
so that $\operatorname{span}(\sigma-\sigma)=\{z=0\}$ after translating $x$ to
the origin.  In these coordinates the condition $\alpha\in K^{p,q}$ says that
every coefficient of $\alpha$ involving only $\dprime y$ and $\ddprime y$
factors vanishes identically on $\RelInt(\sigma)\cap\Omega$, that is, on a
relatively open subset of $\{z=0\}$.  Differentiating such a coefficient in
the directions $y_1,\ldots,y_d$ therefore gives zero there; and the terms of
$\dprime\alpha$ obtained by differentiating in a $z$-direction carry a factor
$\dprime z_i$, which annihilates $N_{\sigma,\mathbf R}$.  Hence
$\dprime\alpha$ evaluates to zero on tangent vectors of $\sigma$ at points of
$\RelInt(\sigma)\cap\Omega$.  If $x\in\sigma\cap\Omega$ lies on the boundary
of $\sigma$, points of $\RelInt(\sigma)\cap\Omega$ approach $x$; evaluating
on a fixed basis of $N_{\sigma,\mathbf R}$ and using continuity gives the
same conclusion at $x$.  This argument is applied separately to every cell,
exactly as required by \eqref{eq:ch10-form-kernel}.  The argument for
$\ddprime$ is identical.  Pullback descends by the chain rule and the fact
that an integer-affine map sends the tangent space of each source cell into
the tangent space of its image cell.  Finally, $J$ is an algebra
automorphism preserving the defining condition.
\end{proof}

The upshot is that $\bigoplus_{p,q}A^{p,q}(\Omega)$ is a bigraded
differential algebra with all the identities of
\Cref{prop:superform-differential-identities}, since those identities hold
before passing to the quotient.

\begin{example}[Smoothness at a tropical vertex]
\label{ex:ch10-smooth-at-vertex}
Consider the tropical line of \Cref{fig:tropical-line-balancing}.  On the
relative interior of each ray, a smooth function is just a smooth
one-variable function.  At the vertex, however, a smooth function in the
present sheaf must be the restriction of one ambient smooth function.
If its outgoing derivatives are $s_1,s_2,s_3$, then
\[
 s_j=\langle \nabla f(0),u_j\rangle,
 \qquad
                         s_1+s_2+s_3=0.
\]
Thus three arbitrary edgewise derivatives do not define a smooth germ at the
vertex.  Piecewise affine functions violating this relation are still
continuous, but their curvature is a delta-current supported at the vertex.
\end{example}

The relation $s_1+s_2+s_3=0$ in that example is the balancing relation
\eqref{eq:ch10-fan-stokes} read on the first jet of an ambient smooth
function.  Indeed, for a balanced one-dimensional fan with rays
$\mathbf R_{\geq0}u_k$ and weights $m_k$, the outgoing derivative of an
ambient smooth function $f$ along the $k$th ray is
$s_k=\langle\nabla f(0),u_k\rangle$, so
\[
 \sum_km_ks_k
 =\Bigl\langle\nabla f(0),\sum_km_ku_k\Bigr\rangle=0 .
\]
For the standard tropical line, this is the only relation among the three
ray vectors.  Hence its ambient smooth germs realize every slope triple in
the hyperplane $s_1+s_2+s_3=0$, as asserted in
\Cref{ex:ch10-smooth-at-vertex}.

For a general star, all linear relations among the ray vectors matter.  Put
$E_0=\operatorname{span}_{\mathbf R}\{u_1,\ldots,u_m\}$.  Elementary
finite-dimensional duality gives
\[
 \bigl\{(\xi(u_k))_{k=1}^m:\xi\in E_0^\vee\bigr\}
 =
 \left\{s\in\mathbf R^m:
   \sum_k\lambda_ks_k=0\ \text{whenever}\
   \sum_k\lambda_ku_k=0\right\}.
\]
Indeed, this is the identity
$\operatorname{im}(U^\vee)=(\ker U)^\perp$ for the linear map
$U:\mathbf R^m\to E_0$ sending the $k$th standard basis vector to $u_k$.
Every covector on $E_0$ extends to the ambient vector space and is the first
jet of an ambient affine function, so the displayed set is exactly the image
of the ambient slope map.
The balancing weights provide one element $(m_k)\in\ker U$, but the kernel
can be larger.  For example, for the cross with ordered rays
$e_1,-e_1,e_2,-e_2$, the ambient slope vectors are precisely
$(a,-a,b,-b)$; the single sum-zero relation is therefore not sufficient.
Thus the smooth sheaf detects the full linear-relation space of the embedded
rays, while the calibration singles out the positive balancing relation used
by integration.

\begin{example}[A form supported on one ray]
\label{ex:ch10-form-one-ray}
On the tropical line $L_2\subset\mathbf R^2$ let $\Omega$ be the complement
of the origin, a disjoint union of three open rays.  By
\Cref{lem:ch10-forms-sheaf}, a form on $\Omega$ is a triple of forms, one on
each ray, with no compatibility condition.  Taking a bump function $\chi$
supported in the relative interior of $\mathbf R_{>0}e_1$ and
$\alpha=\chi\,\dprime x_1\wedge\ddprime x_1$, we get a nonzero element of
$A^{1,1}_c(\Omega)$, which extends by zero to $A^{1,1}_c(L_2)$.  Its integral
against the calibration of weight one is $\int_0^\infty\chi\,\rd s>0$.  This
shows both that $A^{1,1}(L_2)\neq0$ and that the integral of a positive top
form detects the weight of a single ray, as it must if
\Cref{prop:integral-subdivision-independent} is to be sharp.
\end{example}

\section{Tropical spaces and charts}
\label{sec:ch10-tropical-spaces}
\index{tropical space!charts}\index{tropical chart}

Everything so far took place inside one affine space.  A skeleton of a model,
or the tropicalization of an analytic space, is not naturally a subset of a
single $N_{\mathbf R}$: it is covered by pieces that are, and the pieces are
glued by integer-affine maps.  The following definition abstracts that
situation.  \Cref{thm:ch10-forms-presentation-independent} is what makes the
definition usable, because it says that the local data assembled below do not
depend on the chart in which they are read.

\begin{definition}
\label{def:ch10-tropical-space}
An \emph{$n$-dimensional tropical space} is a Hausdorff space $S$ with an
atlas of homeomorphisms $\varphi_i:U_i\to\Omega_i$, where $\Omega_i$ is open
in the support of a pure $n$-dimensional weighted balanced rational
polyhedral complex, such that transition maps are locally restrictions of
integer-affine maps and preserve the weighted cycles.  A superform on $S$ is
a compatible family of superforms in these charts.
\end{definition}

Let us unwind the definition.  An \emph{atlas} is a family
$(U_i,\varphi_i,(\mathscr C_i,m_i))$ with $(U_i)$ an open cover of $S$,
$\varphi_i:U_i\to\Omega_i\subseteq|\mathscr C_i|$ a homeomorphism onto an
open subset, and $(\mathscr C_i,m_i)$ a balanced weighted rational
polyhedral complex of pure dimension $n$ in some $N_{i,\mathbf R}$.  The
transition map on $U_i\cap U_j$ is
\[
 \varphi_{ji}
 =\varphi_j\circ\varphi_i^{-1}:
 \varphi_i(U_i\cap U_j)\longrightarrow\varphi_j(U_i\cap U_j),
\]
and the requirement is that $\varphi_{ji}$ and $\varphi_{ij}$ satisfy
hypotheses \cref{itm:ch10-pres-local} and \cref{itm:ch10-pres-cells} of
\Cref{thm:ch10-forms-presentation-independent}, and in addition that the
bijection of cells there matches weights:
\begin{equation}
\label{eq:ch10-transition-weights}
 m_j(\sigma')=m_i(\sigma)
 \qquad\text{whenever }
 \varphi_{ji}\bigl(\varphi_i(U_i\cap U_j)\cap\sigma\bigr)
 =\varphi_j(U_i\cap U_j)\cap\sigma'
\end{equation}
for a top-dimensional cell $\sigma$ of the refinement of $\mathscr C_i$
supplied by hypothesis \cref{itm:ch10-pres-cells}, with the weights of a
refinement inherited from their carriers as in
\Cref{prop:ch10-balancing-subdivision}.  Two atlases are equivalent if their
union is again an atlas; a tropical space is a space with an equivalence
class of atlases.  Its \emph{level} or \emph{dimension} is $n$.

The following theorem is the verification that
\Cref{def:ch10-tropical-space} defines something: that every structure used
later is transported by transitions, and hence globalizes.
Here a form on a support has one of the three positivity properties if
locally it has an ambient representative in that cone. This is the
representative convention used in \Cref{def:ch13-positivity}.

\begin{theorem}
\label{thm:ch10-transition-transport}
Let $S$ be a tropical space of dimension $n$ with atlas
$(U_i,\varphi_i,(\mathscr C_i,m_i))$.
\begin{enumerate}
\item\label{itm:ch10-tr-forms} The sheaves $U\mapsto A^{p,q}(\varphi_i(U))$
      glue to sheaves $A^{p,q}_S$ whose direct sum is a bigraded algebra
      sheaf on $S$, equipped with
      operators $\dprime,\ddprime$ and an involution $J$ satisfying the
      identities of
      \Cref{prop:superform-differential-identities}.
\item\label{itm:ch10-tr-lattice} For $x\in S$ the tangent lattices of the
      cells through $\varphi_i(x)$, together with their inclusions, are
      independent of $i$ up to canonical isomorphism.  In particular the
      \emph{local weighted star} of $x$ is a well-defined balanced weighted
      fan up to integer-linear isomorphism in its tangent span; unused
      ambient directions are omitted.
\item\label{itm:ch10-tr-calib} The calibrations $\mu_i$ of
      \Cref{def:ch10-calibration} attached to the charts agree on overlaps up
      to the sign of the orientation.  The associated signed density on
      $(n,n)$-forms is chart independent, and is a positive measure when the
      form and the chart weights are positive.
\item\label{itm:ch10-tr-positivity} Positivity, strong positivity, and weak
      positivity of a $(p,p)$-form are chart independent.
\end{enumerate}
\end{theorem}

\begin{proof}
\cref{itm:ch10-tr-forms}
By \Cref{thm:ch10-forms-presentation-independent} applied to $\varphi_{ji}$,
there are isomorphisms
$\varphi_{ji}^*:A^{p,q}(\varphi_j(U_i\cap U_j))\to
A^{p,q}(\varphi_i(U_i\cap U_j))$ compatible with the algebra structure,
$\dprime$, $\ddprime$ and $J$.  Step 4 of that proof shows
$\varphi_{ii}^*=\mathrm{id}$ and, applied to a triple overlap,
$\varphi_{ki}^*=\varphi_{ji}^*\circ\varphi_{kj}^*$, since on
$U_i\cap U_j\cap U_k$ the two composites are induced by integer-affine
extensions restricting to the same map.  This is a descent datum for the
sheaves of \Cref{lem:ch10-forms-sheaf}, so they glue.  The identities of
\Cref{prop:superform-differential-identities} hold in each chart, hence
globally.

\cref{itm:ch10-tr-lattice}
This is Step 5 of the proof of
\Cref{thm:ch10-forms-presentation-independent}: the linear part of a local
integer-affine extension of $\varphi_{ji}$ restricts to an isomorphism
$N_{\sigma}\to N'_{\sigma'}$ for corresponding cells, and the cocycle
condition again holds on triple overlaps.  The isomorphisms are compatible
with the inclusions $N_\tau\subseteq N_\sigma$ for $\tau\prec\sigma$, because
they are restrictions of one linear map.  Consequently the primitive normals
$u_{\sigma/\tau}$ correspond, and by \eqref{eq:ch10-transition-weights} the
weights correspond; the local star of $x$ therefore has a chart-independent
description, and it is balanced because each $(\mathscr C_i,m_i)$ is
balanced and by
\Cref{cor:ch10-balancing-consequences}\,\cref{itm:ch10-bal-local}.

\cref{itm:ch10-tr-calib}
Let $\sigma\leftrightarrow\sigma'$ correspond under $\varphi_{ji}$, with
linear part $A$.  By \cref{itm:ch10-tr-lattice}, $A$ restricts to an
isomorphism $N_\sigma\to N'_{\sigma'}$, so
$\bigwedge^nA$ carries a generator of $\bigwedge^nN_\sigma$ to a generator of
$\bigwedge^nN'_{\sigma'}$, that is,
$\bigwedge^nA(e_\sigma)=\pm e_{\sigma'}$.  Hence
$\bigwedge^nA(\mu_\sigma)=\pm\mu_{\sigma'}$ by
\eqref{eq:ch10-transition-weights}.  The sign is the sign of the determinant
of $A$ in oriented lattice bases and is $+1$ precisely when the two chosen
orientations correspond.  For the measure attached to an $(n,n)$-form, note
that $A$ restricts to an isomorphism of \emph{lattices}
$N_\sigma\to N'_{\sigma'}$, so its determinant computed in lattice bases is
$\pm1$; the pullback of the interlaced volume form $\tau_n$ therefore
acquires the factor $\det(A)^2=1$ from the two labelled exterior blocks, the
lattice-normalized Lebesgue measures correspond, and the weight is
unchanged.  The resulting signed density is thus chart independent, and it
is a positive measure when both the form and the weights are positive. This is
\Cref{prop:appG-calibration-tests} in the present notation, and it is the
precise sense in which an orientation is bookkeeping while a weight is not.

\cref{itm:ch10-tr-positivity}
All three positivity cones are stable under pullback by an affine map, by the
corresponding statement in \Cref{chap:superforms}; and by Step 3 of the proof
of \Cref{thm:ch10-forms-presentation-independent} the pullback on
$A^{p,q}(\Omega)$ does not depend on the affine extension chosen.  Since
$\varphi_{ji}^*$ and $\varphi_{ij}^*$ are mutually inverse, each cone is
carried onto the corresponding cone.
\end{proof}

\begin{definition}
\label{def:ch10-tropical-morphism}
A \emph{morphism} of tropical spaces $f:S\to T$ is a continuous map such
that, for every $x\in S$, there are charts $\varphi_i$ near $x$ and $\psi_j$
near $f(x)$ for which $\psi_j\circ f\circ\varphi_i^{-1}$ is the restriction
of an integer-affine map.  A morphism is \emph{compact} if it is proper, and
\emph{generically finite on a cell} if its linear part is
injective on the tangent space of that cell.
\end{definition}

After refining cells so that each source cell maps into a target cell,
pullback along a morphism is defined by
\Cref{prop:ch10-forms-quotient-operations}. Independence of the affine
extension and of the charts follows as in Steps 2--3 of
\Cref{thm:ch10-forms-presentation-independent}.
Pushforward of cycles requires the lattice indices of
\Cref{sec:ch10-pushforward}, and it is not defined for arbitrary morphisms.

\begin{example}[The tropical line as a tropical space]
\label{ex:ch10-line-as-space}
The support $L_2\subset\mathbf R^2$ of \Cref{ex:ch10-standard-line} is a
tropical space with a one-chart atlas. An alternative atlas starts with
the three open rays, each mapped to an open subset of $\mathbf R$
with weight one: for the ray $\mathbf R_{>0}e_1$, use $x\mapsto x_1$; for
$\mathbf R_{>0}e_2$, use $x\mapsto x_2$; for $\mathbf R_{>0}e_0$, use
$x\mapsto-x_1$.  These three charts cover $L_2\setminus\{0\}$ only, so a
fourth chart, namely the identity on a neighbourhood of the vertex, is
needed.  The transition between the fourth chart and the first is
$x\mapsto x_1$ on $\{x_2=0,\,x_1>0\}$, an integer-linear map carrying the
tangent lattice $\mathbf Z e_1$ isomorphically onto $\mathbf Z$ and the
weight one to the weight one.  Thus the two atlases are equivalent, as
\Cref{thm:ch10-transition-transport} requires.
\end{example}

\begin{example}[A tropical space that is not a polyhedral subset]
\label{ex:ch10-circle}
Let $\ell>0$ and let $S=\mathbf R/\ell\mathbf Z$, covered by the images of
three open intervals of length $2\ell/3$; each is mapped isomorphically to an
open subset of $\mathbf R$, carrying the standard lattice $\mathbf Z$ and
weight one, and the transitions are translations, which are integer affine
with integral linear part $+1$.  So $S$ is a one-dimensional tropical space.

It is not homeomorphic to the support of any balanced weighted complex with
positive weights, in any $N_{\mathbf R}$.  We prove more: such a support is
never nonempty and compact.

Suppose $(\mathscr C,m)$ is balanced, pure of dimension one, with $m_\sigma>0$
for every facet, and with $|\mathscr C|$ compact and nonempty.  Local
finiteness and compactness give finitely many cells.  Choose $a\in M$ such
that $\langle a,u_\sigma\rangle\neq0$ for the primitive direction of every
edge $\sigma$; this is possible because only finitely many proper hyperplanes
of $M_{\mathbf R}$ are excluded.  The linear functional
$x\mapsto\langle a,x\rangle$ attains its maximum on $|\mathscr C|$.  It is
nonconstant on every edge, so a maximum occurs at a vertex $x$.  For every
edge $\sigma$ incident to $x$, let $u_{\sigma/x}$ be its primitive direction
pointing away from $x$.  Maximality and the generic choice of $a$ give
$\langle a,u_{\sigma/x}\rangle<0$.  Pairing the balancing equation at $x$
with $a$ now yields
\[
 0=\left\langle a,
      \sum_{\sigma\ni x}m_\sigma u_{\sigma/x}\right\rangle
  =\sum_{\sigma\ni x}m_\sigma
     \langle a,u_{\sigma/x}\rangle<0,
\]
a contradiction.  Hence no such complex exists, and in particular the
circle is not such a support.

The hypothesis of positive weights is used, and is not decorative: with
weights of mixed sign a vertex of the convex hull of the support can satisfy
the balancing equation, as the primitive vectors $(-1,0)$, $(-1,1)$,
$(-1,-1)$ with weights $-2,1,1$ show.  Tropicalizations of analytic spaces
carry positive multiplicities
(\Cref{thm:appG-analytic-calibration}), so the effective case is the
relevant one.

The circle $\mathbf R/\ell\mathbf Z$ is exactly the skeleton of a Tate curve
computed in \Cref{ex:ch08-tate-curve}, with $\ell$ the valuation of the Tate
parameter.  Thus the smallest interesting non-Archimedean example already
forces the abstract definition: this Tate skeleton is a tropical space
that is not an effective balanced tropical subset of $\mathbf R^r$.
\end{example}

This economical definition suffices for the smooth calculus. CLD's parlinear
and piecewise-linear spaces retain additional sheaf data needed for maximum
generality. We invoke that formalism only when a theorem genuinely requires
it.  The relation is the following: a parlinear space carries, in addition to
the topological and chart data above, a sheaf of parlinear functions, and the
tropical charts are required to be sections of it; the balancing condition is
replaced by the harmony condition on the calibration.  For the theorems of
\Cref{chap:tropical-integration,chap:cld-forms} the two formalisms give the
same integrals, and we use whichever is lighter.  The maximally general
version, together with the vertebrate and chorded variants needed for spaces
with boundary, is \cite[\S\S3--6]{CLD25}; a précis is
\Cref{sec:appG-polyhedra}.

\section{Calibrations, orientations, and multiplicities}
\label{sec:ch10-calibrations}
\index{calibration}\index{orientation!of a polyhedral cell}

The weights of a balanced complex were introduced as integers attached to
facets.  To transport oriented top-degree data one also chooses primitive
generators of the top exterior tangent lattices.  A calibration records the
product of the chosen generator and the weight, so that this data can be
transported by a chart change all at once, as in
\Cref{thm:ch10-transition-transport}\,\cref{itm:ch10-tr-calib}.  We introduce
it here, immediately before the pushforward and intersection theory that
consume it, and immediately before the integration of
\Cref{chap:tropical-integration} where its purpose is visible.

\begin{definition}
\label{def:ch10-calibration}
An orientation of an $n$-cell is a choice of one of the two primitive
generators $e_\sigma$ of $\bigwedge^nN_\sigma$.  After making this choice, a
\emph{calibration} assigns to each maximal cell the multivector
\[
                              \mu_\sigma=m_\sigma e_\sigma,
\]
compatibly with subdivision.  For a positive cycle its magnitude is the
positive weight; for a signed cycle its coefficient relative to the chosen
orientation is the signed weight.
\end{definition}

A calibration is an oriented multiplicity in the top exterior tangent
lattice.  Its factorization into an orientation and a signed cycle weight is
not canonical unless one of those two data is fixed, as the next proposition
makes precise.  Reversing orientation changes integrals of ordinary
$n$-forms; it does not change the density attached to an
$(n,n)$-form because both labelled copies reverse.

\begin{proposition}
\label{prop:ch10-calibration-basics}
Let $\mathscr C$ be a pure $n$-dimensional rational polyhedral complex in
$N_{\mathbf R}$.
\begin{enumerate}
\item\label{itm:ch10-cal-bijection} A pair consisting of a weighting $m$ and
      a choice of orientations determines the family
      $\mu_\sigma=m_\sigma e_\sigma$ of nonzero elements of
      $\bigwedge^nN_\sigma$.  On each facet every such element arises from
      exactly two pairs, differing by the simultaneous sign change
      $(m_\sigma,e_\sigma)\mapsto(-m_\sigma,-e_\sigma)$; these choices are
      independent from facet to facet.  In particular the calibration
      determines $|m_\sigma|$ but not the sign of $m_\sigma$; that sign is
      the datum which distinguishes a cycle from its negative.
\item\label{itm:ch10-cal-subdivision} If $\mathscr C'$ subdivides
      $\mathscr C$ and each facet $\sigma'$ of $\mathscr C'$ receives
      $\mu_{\sigma'}=\mu_\sigma$ for its carrier $\sigma$, using the
      inherited orientation $e_{\sigma'}=e_\sigma$ and weight
      $m_{\sigma'}=m_\sigma$, then the
      assignment is a calibration on $\mathscr C'$, and the two calibrations
      define the same integral of any compactly supported $(n,n)$-form.
\item\label{itm:ch10-cal-balancing} With the incidence numbers
      \eqref{eq:ch10-incidence} computed from the same orientations, the
      complex is balanced if and only if
      $\sum_{\sigma\supset\tau}[\sigma:\tau]\mu_\sigma=0$ in
      $\bigwedge^nN$ for every codimension-one cell $\tau$.  Each product
      $[\sigma:\tau]\mu_\sigma$ is independent of the orientation
      $e_\sigma$, the weighting $m$ being held fixed.
\end{enumerate}
\end{proposition}

\begin{proof}
\cref{itm:ch10-cal-bijection} By \Cref{lem:ch10-tangent-lattice}\,(2) the
group $\bigwedge^nN_\sigma$ is free of rank one, with exactly two
generators $\pm e_\sigma$; a nonzero element of it is therefore
$m_\sigma e_\sigma$ for a unique nonzero integer $m_\sigma$ once a
generator has been chosen, and the only ambiguity is replacing
$(e_\sigma,m_\sigma)$ by $(-e_\sigma,-m_\sigma)$.  This choice is made
independently on each facet.  Hence $|m_\sigma|$ is determined by
$\mu_\sigma$ and its sign is not.

\cref{itm:ch10-cal-subdivision} The assignment makes sense because
$N_{\sigma'}=N_\sigma$ by \Cref{lem:ch10-subdivision-lattices}, so
$\mu_\sigma$ lies in $\bigwedge^nN_{\sigma'}$.  The integral statement is
\Cref{prop:integral-subdivision-independent}: the relative interiors of the
$\sigma'$ with carrier $\sigma$ partition $\RelInt(\sigma)$ up to a set of
measure zero, and the lattice-normalized measures agree.

\cref{itm:ch10-cal-balancing} This is
\Cref{thm:ch10-balancing-boundary} rewritten with
$\mu_\sigma=m_\sigma e_\sigma$; the independence of the products from the
orientation is Step~4 of that proof, in which $m$ is held fixed and both
$[\sigma:\tau]$ and $\mu_\sigma$ change sign together.
\end{proof}

\subsection{The sign conventions, tested}
\label{sec:ch10-sign-test}

Sign conventions are only trustworthy after being run on an example in which
every choice is made explicitly.  We do that now, on the smallest example in
which orientation, weight, and lattice all appear.

\begin{example}[Orientation reversal on a segment]
\label{ex:ch10-sign-test-segment}
Work in $N=\mathbf Z$, $N_{\mathbf R}=\mathbf R$, on the cell
$\sigma=[0,2]$ with weight $m_\sigma=3$.  Choose the orientation
$e_\sigma=+1\in\bigwedge^1\mathbf Z=\mathbf Z$, so that
$\mu_\sigma=3$.  Let $\alpha=f(x)\,\dprime x\wedge\ddprime x$ with $f\geq0$
smooth.  With the conventions of
\Cref{chap:tropical-integration},
\[
 \int_{(\sigma,\mu_\sigma)}\alpha
 =3\int_0^2f(x)\,\rd x\geq0 .
\]
Now reverse the orientation, keeping the weight: $e_\sigma'=-1$,
$m_\sigma=3$, $\mu_\sigma'=-3$.  The integral is unchanged.  Indeed, in the
reversed coordinate $y=-x$ one has
$\dprime y\wedge\ddprime y=(-\dprime x)\wedge(-\ddprime x)
=\dprime x\wedge\ddprime x$, so the interlaced volume form $\tau_1$ is the
same and the lattice-normalized Lebesgue measure is the same.  This is
\Cref{prop:appG-calibration-tests}\,(3): an $(n,n)$-form does not see an
orientation, because the two labelled exterior blocks reverse together.

Next let $\eta=g(x)\,\ddprime x$ be a form of bidegree $(0,1)$.  The
following identity is the weight-one cell integral; the calibrated identity
multiplies both sides by $3$. The boundary traces
\eqref{eq:superform-boundary-traces} with $n=1$ read
$\operatorname{tr}'_{\sigma,\tau}(\eta)=(\iota''_{n_{\sigma,\tau}}\eta)|_\tau$,
the prefactor $(-1)^{n-1}$ being $+1$ in this dimension, with
$n_{\sigma,\tau}$ the outward normal, that is, $+1$ at $x=2$ and $-1$ at
$x=0$.  Hence
\[
 \int_\sigma\dprime\eta
 =\int_0^2g'(x)\,\rd x
 =g(2)-g(0)
 =\operatorname{tr}'_{\sigma,\{2\}}(\eta)
  +\operatorname{tr}'_{\sigma,\{0\}}(\eta),
\]
in agreement with \Cref{thm:polyhedral-stokes}.  Reversing the orientation of
$\sigma$ changes neither side.  The traces are computed from the outward
primitive normals $n_{\sigma,\tau}$, which are geometric and independent of
$\operatorname{or}(\sigma)$; and both integrals are unsigned, by
\Cref{lem:ch11-orientation-drops-out}.  What does flip is the induced
orientation of each endpoint, and that enters neither integral.  This is
the assertion that orientation reversal is invisible on both sides of the
Stokes pairing --- as it must be, the left-hand side being the integral of
an $(n,n)$-form, shown orientation-blind two paragraphs above.
\end{example}

\begin{example}[The three orientations of the tropical line]
\label{ex:ch10-sign-test-line}
On $L_2\subset\mathbf R^2$ take the three rays with weights $1$ and the
orientations $e_1$, $e_2$, $e_0=-e_1-e_2$ pointing away from the vertex, so
that $\mu_k=e_k$ for $k=0,1,2$.  The incidence numbers
\eqref{eq:ch10-incidence} at the vertex $\tau=\{0\}$ are all $-1$: with
$e_\tau=1\in\bigwedge^0N$ and $\widetilde u_{\sigma_k/\tau}=e_k=e_{\sigma_k}$,
the defining relation
$\widetilde u_{\sigma_k/\tau}\wedge e_\tau=-[\sigma_k:\tau]e_{\sigma_k}$
gives $[\sigma_k:\tau]=-1$.  Hence
\[
 \sum_{k=0}^2[\sigma_k:\tau]\mu_k
 =-(e_0+e_1+e_2)=0,
\]
confirming \Cref{prop:ch10-calibration-basics}\,\cref{itm:ch10-cal-balancing}.

Now reverse the orientation of the ray $\mathbf R_{\geq0}e_0$ only, keeping
the weight $m_0=1$: the new data are $e_{\sigma_0}'=-e_0$,
$\mu_0'=-e_0$, and the new incidence number is
$[\sigma_0:\tau]'=+1$, because
$\widetilde u_{\sigma_0/\tau}\wedge e_\tau=e_0
=-[\sigma_0:\tau]'\,(-e_0)$.  The product is
$[\sigma_0:\tau]'\mu_0'=-e_0$, exactly what it was before; the calibrated
balancing sum is therefore unchanged, as
\Cref{prop:ch10-calibration-basics}\,\cref{itm:ch10-cal-balancing}
predicts.  The uncalibrated sum $\sum_km_ku_{\sigma_k/\tau}$ never involved
the orientations at all, and is likewise unchanged.

Contrast this with changing the \emph{weight} of that ray from $1$ to $-1$,
keeping the orientation $e_0$.  Now $\mu_0=-e_0$ and
$[\sigma_0:\tau]=-1$, so the calibrated sum becomes
$e_0-e_1-e_2=2e_0\neq0$, while the uncalibrated sum becomes
$-e_0+e_1+e_2=-2e_0\neq0$.  Both detect the failure, as they must: the
weighting $(1,1,-1)$ is a different weighted complex and is not a tropical
cycle, since it is not balanced.

The moral is the one recorded in the warning below.  Reversing an orientation
is invisible to every geometric quantity---the calibrated balancing equation,
the integral of an $(n,n)$-form, the measure attached to a positive
form---whereas changing the sign of a weight changes the cycle.  When a
computation produces an unexpected factor of $-1$, the first question is
which of the two has happened.
\end{example}

\begin{warning}{Three different signs}
Three signs occur in this chapter and must be kept apart: the sign of a
weight $m_\sigma$, which distinguishes a cycle from its negative; the sign of
an orientation $e_\sigma$, which is a bookkeeping choice; and the sign of an
incidence number $[\sigma:\tau]$, which is determined by the orientations.
Only the first has geometric content.  Flipping $e_\sigma$ flips both
$\mu_\sigma$ and $[\sigma:\tau]$, so it leaves the products
$[\sigma:\tau]\mu_\sigma$---and hence the balancing equation, the boundary
\eqref{eq:ch10-boundary}, and all integrals of $(n,n)$-forms---unchanged.
Flipping $m_\sigma$ does not.  Note also that the calibration $\mu_\sigma$
alone determines $|m_\sigma|$ but not its sign, so a calibration is not by
itself enough to test balancing; the weighting must be specified as well.
For an effective cycle this ambiguity disappears once one agrees that all
weights are positive.
\end{warning}

\section{Integer-affine pushforward}
\label{sec:ch10-pushforward}
\index{pushforward!of tropical cycles}\index{lattice index}

Pushforward is the operation in which lattice indices first appear, and it is
the operation most often mishandled.  We isolate the index algebra in a
lemma, prove the pushforward statement, and then test it on examples where
the answer can be checked independently.

\begin{lemma}
\label{lem:ch10-index}
Let $N$ be a free abelian group of finite rank and let $A,B,C\subseteq N$ be
subgroups of finite index in their saturations.
\begin{enumerate}
\item\label{itm:ch10-idx-tower} If $A\subseteq B\subseteq C$ have equal
      ranks, then $[C:A]=[C:B][B:A]$.
\item\label{itm:ch10-idx-quotient} If $A$ and $B$ are saturated with
      $A+B$ of finite index in $N$, then the natural map
      $N/(A+B)\to(N/A)/\overline B$ is an isomorphism, where $\overline B$ is
      the image of $B$; in particular
      $[N:A+B]=[N/A:\overline B]$.
\item\label{itm:ch10-idx-kernel-cokernel} If $A,B,C\subseteq N$ are saturated
      and $A+B$, $B+C$, $(A\cap B)+C$ all have finite index in $N$, then
      \[
       [N:(A\cap B)+C]
       =[N:B+C]\,\bigl[B:(A\cap B)+(B\cap C)\bigr].
      \]
\item\label{itm:ch10-idx-diagonal} Let $\Delta=\{(a,a):a\in N\}\subseteq
      N\oplus N$ and let $A,B\subseteq N$ be saturated with $A+B$ of finite
      index in $N$. Then
      \[
       \bigl[N\oplus N:(A\oplus B)+\Delta\bigr]=[N:A+B].
      \]
\end{enumerate}
\end{lemma}

\begin{proof}
\cref{itm:ch10-idx-tower} is the multiplicativity of indices in a tower of
finite-index subgroups.

\cref{itm:ch10-idx-quotient} The composite $N\to N/A\to(N/A)/\overline B$ is
surjective with kernel $A+B$, giving the isomorphism, and the index
statement is the resulting equality of orders.

\cref{itm:ch10-idx-kernel-cokernel} Consider the homomorphism
\[
 \theta:B\longrightarrow N/\bigl((A\cap B)+C\bigr)
\]
induced by the inclusion $B\subseteq N$.  Its kernel is
$B\cap((A\cap B)+C)$, which contains $(A\cap B)+(B\cap C)$; conversely, if
$b=a'+c$ with $a'\in A\cap B$ and $c\in C$, then $c=b-a'\in B\cap C$, so the
kernel is exactly $(A\cap B)+(B\cap C)$.  Its cokernel is
$N/\bigl((A\cap B)+C+B\bigr)=N/(B+C)$, since $A\cap B\subseteq B$.  Hence
there is an exact sequence
\[
 0\to(A\cap B)+(B\cap C)\to B\to N/\bigl((A\cap B)+C\bigr)\to N/(B+C)\to0
\]
of abelian groups in which the first and last quotients are finite; comparing
orders in the induced short exact sequence
\[
 0\to B/\bigl((A\cap B)+(B\cap C)\bigr)\to
 N/\bigl((A\cap B)+C\bigr)\to N/(B+C)\to0
\]
gives the asserted identity.

\cref{itm:ch10-idx-diagonal} The map $\delta:N\oplus N\to N$,
$\delta(a,b)=a-b$, is surjective with kernel $\Delta$, and
$\delta(A\oplus B)=A+B$.  Hence $\delta$ induces an isomorphism
$(N\oplus N)/\bigl((A\oplus B)+\Delta\bigr)\cong N/(A+B)$, and the orders
agree.
\end{proof}

Let $F:N_{\mathbf R}\to N'_{\mathbf R}$ be integer affine.  Suppose that
$F$ is injective on the span of every maximal cell of a pure
$n$-dimensional weighted complex $C$. Assume in addition that $C$ is finite,
or that $F:|C|\to N'_{\mathbf R}$ is proper. After subdivision, the images
of the maximal cells form a complex. In the proper case they are locally
finite: a compact target neighbourhood has compact inverse image, which
meets only finitely many source cells. For an image cell $\sigma'$ define
\[
 m_{\sigma'}(F_*C)
 =\sum_{\substack{\sigma\in C_n\\F(\sigma)=\sigma'}}
   m_\sigma\,[N_{\sigma'}:F(N_\sigma)].
\]
The equality $F(\sigma)=\sigma'$ here means equality after restricting to the
relative interiors of suitably subdivided cells.

The extra hypothesis matters even for balanced complexes. The horizontal
lines $\mathbf R\times\{j\}$, $j\in\mathbf Z$, with weight one form a
locally finite balanced complex. Projection onto the first coordinate is
injective on each tangent line, but would sum infinitely many unit weights
on the image line. This projection is not proper on the support.

\begin{proposition}
\label{prop:tropical-pushforward-balanced}
Under these finiteness or properness hypotheses, if $C$ is balanced then
$F_*C$ is balanced. The construction is unchanged
by subdivision and is functorial for compositions satisfying the same
dimension and finiteness or properness hypotheses.
\end{proposition}
\begin{proof}
Fix a codimension-one cell $\tau'$ of $F_*C$ and a facet
$\sigma'\supset\tau'$.  Let $\tau\subseteq F^{-1}(\tau')$ be a
codimension-one cell of $C$ with $F(\tau)=\tau'$, and let
$\sigma\supset\tau$ be a facet with $F(\sigma)=\sigma'$.  Because $F$ is
injective on $N_{\sigma,\mathbf R}$, the induced map
\[
 \bar F:N_\sigma/N_\tau\longrightarrow N'_{\sigma'}/N'_{\tau'}
\]
is injective between free groups of rank one; hence
$\bar F(u_{\sigma/\tau})=\lambda_{\sigma}\,u_{\sigma'/\tau'}$ for a positive
integer $\lambda_\sigma$, the sign being positive because $F$ maps the
$\sigma$ side of $\tau$ into the $\sigma'$ side of $\tau'$.  Applying
\Cref{lem:ch10-index}\,\cref{itm:ch10-idx-tower} to the inclusions
$F(N_\sigma)\subseteq F(N_\sigma)+N'_{\tau'}\subseteq N'_{\sigma'}$ and
comparing the induced rank-one quotients gives
\[
 [N'_{\sigma'}:F(N_\sigma)]\,\lambda_\sigma^{-1}
 =[N'_{\tau'}:F(N_\tau)],
\]
that is, $[N'_{\sigma'}:F(N_\sigma)]=\lambda_\sigma[N'_{\tau'}:F(N_\tau)]$.
Now apply $\bar F$ to the source balancing equation at $\tau$ and multiply by
$[N'_{\tau'}:F(N_\tau)]$:
\[
 0=[N'_{\tau'}:F(N_\tau)]\sum_{\sigma\supset\tau}m_\sigma
     \bar F(u_{\sigma/\tau})
  =\sum_{\sigma\supset\tau}m_\sigma
     [N'_{\sigma'(\sigma)}:F(N_\sigma)]
       u_{\sigma'(\sigma)/\tau'},
\]
where $\sigma'(\sigma)=F(\sigma)$; summing over the finitely many $\tau$
lying over $\tau'$ regroups the left-hand side into the balancing sum of
$F_*C$ at $\tau'$, whose coefficients are exactly the pushforward weights.

For subdivision invariance, pass to a common subdivision on which all images
are cells.  A refined facet has the same tangent lattice and weight as its
carrier by \Cref{lem:ch10-subdivision-lattices}; its relative interior lies
in exactly one refined image facet.  Grouping refined facets by their old
carriers therefore reproduces every term in the displayed pushforward
formula exactly once.

Finally let $G:N'_{\mathbf R}\to N''_{\mathbf R}$ satisfy the same dimension
hypothesis.  For a chain of image facets
$\sigma\mapsto\sigma'\mapsto\sigma''$, the restriction of $G$ to
$N'_{\sigma'}$ is injective, so it preserves the finite index of
$F(N_\sigma)$ in $N'_{\sigma'}$.  The tower formula gives
\[
 [N''_{\sigma''}:G(F(N_\sigma))]
 =[N''_{\sigma''}:G(N'_{\sigma'})]\,
   [N'_{\sigma'}:F(N_\sigma)].
\]
Summing first over the facets mapping to $\sigma'$ and then over the
$\sigma'$ mapping to $\sigma''$ is therefore the same as summing directly
for $G\circ F$.  This proves functoriality.
\end{proof}

\begin{example}[A degree-two lattice map]
\label{ex:ch10-degree-two}
Let $F:\mathbf R\to\mathbf R$ be $t\mapsto2t$, with source and target lattice
$\mathbf Z$.  The image ray has the same underlying set, but
$[\mathbf Z:2\mathbf Z]=2$.  A source ray of weight one therefore pushes
forward to a ray of weight two.  This index is the tropical counterpart of a
degree in the analytic projection formula.
\end{example}

\begin{example}[Projecting the tropical line]
\label{ex:ch10-project-line}
Let $L_2\subset\mathbf R^2$ be the standard tropical line with weights one
and let $F:\mathbf R^2\to\mathbf R$ be $F(x_1,x_2)=x_1$.  The map $F$ is
injective on the tangent line of the ray $\mathbf R_{\geq0}e_1$ and on that
of $\mathbf R_{\geq0}e_0$, but it kills the tangent line of
$\mathbf R_{\geq0}e_2$.  The hypothesis of
\Cref{prop:tropical-pushforward-balanced} therefore fails, and indeed the
naive image of $L_2$ is all of $\mathbf R$ with the ray
$\mathbf R_{\geq0}e_2$ collapsing to the origin.

To illustrate the stated construction, use instead the balanced one-cycle
$Z=\mathbf R_{\geq0}e_1\cup\mathbf R_{\geq0}(-e_1)$, with weights one; here
$F$ is injective on both tangent lines, $[\mathbf Z:F(\mathbf Ze_1)]=1$, and
$F_*Z=\mathbf R$ with weight one, which is balanced.  If instead one pushes
forward $Z'=\mathbf R_{\geq0}(2e_1)\cup\mathbf R_{\geq0}(-2e_1)$---the same
set, but presented with the nonsaturated lattice $\mathbf Z(2e_1)$, which is
\emph{not} allowed by our conventions---one obtains weight two.  This is the
error against which the warning after \Cref{def:appG-affine-conventions}
guards: primitivity is computed in the saturated lattice.
\end{example}

\begin{example}[A tropical plane projected to a plane]
\label{ex:ch10-project-plane}
Let $H\subset\mathbf R^3$ be the standard tropical plane of
\Cref{ex:ch10-standard-plane} and let $F:\mathbf R^3\to\mathbf R^2$ be the
projection $(x_1,x_2,x_3)\mapsto(x_1,x_2)$.  The linear part is injective on
the tangent plane of $\operatorname{cone}(e_1,e_2)$, of
$\operatorname{cone}(e_1,e_0)$ and of $\operatorname{cone}(e_2,e_0)$, and it
kills $e_3$, so it is not injective on the three cones containing $e_3$.  The
hypothesis of \Cref{prop:tropical-pushforward-balanced} fails, so its
construction does not apply. More general tropical pushforward rules may
discard collapsed facets, but require their own statement and proof.
Nor does the failure force every top-degree integral to diverge: a pulled
back $(2,2)$-form restricts to zero on a facet whose image has dimension one.
Properness is a separate condition in the analytic projection formula.
\end{example}

\section{Stable intersections}
\label{sec:ch10-stable-intersection}
\index{stable intersection}\index{fan displacement rule}

Two tropical cycles rarely meet in the expected dimension.  The tropical line
$L_2$ and its translate by a vector along the ray $\mathbf R_{\geq0}e_0$
share a whole ray, whereas the expected dimension of an intersection of two
curves in the plane is zero.  Classical intersection theory solves this by
moving one cycle; the tropical theory does the same, and the surprise is that
the answer is independent of how one moves, provided the cycles are balanced.
That is the content of \Cref{thm:tropical-stable-intersection}; its
wall-crossing calculation identifies the role of balancing in the moving
construction.

Following the organizing principle of \Cref{part:cld}, we do the computations
first.  \Cref{sec:ch10-transverse} defines the transverse multiplicity, which
requires no displacement.  \Cref{sec:ch10-three-computations} works out three
intersections completely by hand, two in $\mathbf R^2$ and one in
$\mathbf R^3$, and verifies in each case that the count does not depend on
the displacement direction.  Only then, in
\Cref{sec:ch10-general-stable}, is the general construction stated and
proved.

\subsection{Transverse intersections}
\label{sec:ch10-transverse}

Let $X$ and $Y$ be tropical cycles of pure dimensions $k$ and $l$ in
$N_{\mathbf R}$, initially represented by finite complexes.  They meet
\emph{transversely} if every facet of $X\cap Y$ has dimension
$k+l-r$, is uniquely an intersection $\sigma\cap\sigma'$ of facets of
$X$ and $Y$, and has relative interior contained in both
$\RelInt(\sigma)$ and $\RelInt(\sigma')$, with
$N_{\sigma,\mathbf R}+N_{\sigma',\mathbf R}=N_{\mathbf R}$. Its weight is
\[
 \omega(\sigma\cap\sigma')
 =\omega_X(\sigma)\omega_Y(\sigma')
   [N:N_\sigma+N_{\sigma'}].
\]

The index $[N:N_\sigma+N_{\sigma'}]$ is finite exactly when
$N_{\sigma,\mathbf R}+N_{\sigma',\mathbf R}=N_{\mathbf R}$, which is the
linear-algebra form of transversality; and it equals one exactly when
$N_\sigma$ and $N_{\sigma'}$ together generate $N$.  The next example shows
that the second condition is strictly stronger, so the index cannot be
dropped.

\begin{example}[The lattice-index test]
\label{ex:ch10-lattice-index-test}
The coordinate axes in $\mathbf R^2$ meet with multiplicity one. The lines
spanned by $(1,0)$ and $(1,2)$ meet transversely with multiplicity
$|\det((1,0),(1,2))|=2$. Any definition that gives one in the second example
has forgotten the lattice calibration.
\end{example}

In $\mathbf R^r$, when $\dim\sigma+\dim\sigma'=r$, the index is computed as an
absolute determinant: if $a_1,\ldots,a_k$ is a lattice basis of $N_\sigma$ and
$b_1,\ldots,b_{r-k}$ one of $N_{\sigma'}$, then
\begin{equation}
\label{eq:ch10-index-determinant}
 [N:N_\sigma+N_{\sigma'}]
 =\bigl|\det(a_1,\ldots,a_k,b_1,\ldots,b_{r-k})\bigr|,
\end{equation}
the determinant being taken in a lattice basis of $N$.  Indeed, the columns
listed generate $N_\sigma+N_{\sigma'}$, and the index of a finite-index
subgroup of $\mathbf Z^r$ is the absolute value of the determinant of a
generating matrix.  This is the formula used in every computation below.

\subsection{Three computations}
\label{sec:ch10-three-computations}

\begin{example}[Two tropical lines in general position]
\label{ex:ch10-two-lines-generic}
Let $L=L_2$ be the standard tropical line in $\mathbf R^2$, with rays
\[
 R_1=\mathbf R_{\geq0}(1,0),\quad
 R_2=\mathbf R_{\geq0}(0,1),\quad
 R_0=\mathbf R_{\geq0}(-1,-1),
\]
all of weight one, and let $L'=L+(3,1)$, the translate with vertex $(3,1)$
and rays $R_i'=R_i+(3,1)$.  We claim that $L$ and $L'$ meet transversely in
the single point $(2,0)$, with multiplicity one.

Write the rays as sets:
\[
\begin{aligned}
 R_1&=\{(t,0):t\geq0\}, & R_2&=\{(0,t):t\geq0\}, &
 R_0&=\{(s,s):s\leq0\},\\
 R_1'&=\{(3+t,1):t\geq0\}, & R_2'&=\{(3,1+t):t\geq0\}, &
 R_0'&=\{(s,s-2):s\leq3\}.
\end{aligned}
\]
We test the nine pairs.  The pairs $(R_1,R_1')$ and $(R_2,R_2')$ have
parallel tangent lines and disjoint supports.  For $(R_1,R_2')$ one needs
$x_1=3$ and $x_2=0$, but $R_2'$ has $x_2\geq1$.  For $(R_2,R_1')$ one needs
$x_1=0$ and $x_1\geq3$.  For $(R_0,R_0')$ one needs $s=s-2$.  The pair
$(R_1,R_0')$ gives $x_2=0$, hence $s=2$ and the point $(2,0)$, which lies in
$R_1$ because $2\geq0$ and in $R_0'$ because $2\leq3$.  The pair
$(R_2,R_0')$ gives $x_1=0$, hence $s=0$ and the point $(0,-2)$, which is not
in $R_2$.  The pair $(R_0,R_1')$ gives $x_2=1$ and $x_1=x_2$, that is
$(1,1)$, which is not in $R_0$; the pair $(R_0,R_2')$ gives $x_1=3$, hence
the point $(3,3)$, again not in $R_0$.

So $L\cap L'=\{(2,0)\}$, a single point, of the expected dimension
$1+1-2=0$, arising from exactly one pair of facets; the intersection is
transverse.  Its multiplicity is
\[
 1\cdot1\cdot\bigl|\det\bigl((1,0),(-1,-1)\bigr)\bigr|
 =\bigl|1\cdot(-1)-0\cdot(-1)\bigr|=1,
\]
using \eqref{eq:ch10-index-determinant} with the primitive generators
$(1,0)$ of $N_{R_1}$ and $(-1,-1)$ of $N_{R_0'}$.  The total multiplicity is
therefore $1$, which is the tropical Bézout count for two lines.
\end{example}

\begin{example}[Two tropical lines in degenerate position]
\label{ex:ch10-two-lines-degenerate}
Keep $L$ as above and let $a>0$, $L'=L+(a,a)$.  Now
\[
 R_0=\{(s,s):s\leq0\}\subset R_0'=\{(s,s):s\leq a\},
\]
so $L\cap L'\supseteq R_0$ is one-dimensional: the intersection is not
transverse, and the naive count is not defined.  We displace.

Let $v\in\mathbf R^2$ and $\varepsilon>0$, and intersect $L$ with
$L'+\varepsilon v$, the tropical line with vertex
$p_\varepsilon=(a+\varepsilon v_1,\,a+\varepsilon v_2)$.  We compute for
three choices of $v$.

\emph{The direction $v=(1,0)$.}  The vertex is $(a+\varepsilon,a)$ and the
three rays are $\{(a+\varepsilon+t,a)\}$, $\{(a+\varepsilon,a+t)\}$, and
$\{(s+\varepsilon,s):s\leq a\}$.  Intersecting with $R_1=\{(t,0)\}$: the
first ray has $x_2=a>0$; the second has $x_1=a+\varepsilon>0$ and $x_2\geq
a>0$; the third gives $s=0$, hence the point $(\varepsilon,0)$, which lies in
$R_1$.  Intersecting with $R_2=\{(0,t)\}$: the third ray gives
$s=-\varepsilon$, hence $x_2=-\varepsilon<0$, not in $R_2$; the first two
rays have $x_1>0$.  Intersecting with $R_0=\{(s,s):s\leq0\}$: the first ray
forces $s=a>0$; the second forces $s=a+\varepsilon>0$; the third has
$x_1-x_2=\varepsilon\neq0$.  So the only intersection point is
$(\varepsilon,0)$, transverse, with multiplicity
$|\det((1,0),(-1,-1))|=1$.  Total: $1$.

\emph{The direction $v=(0,1)$.}  Symmetrically, the vertex is
$(a,a+\varepsilon)$, the third ray is $\{(s,s+\varepsilon):s\leq a\}$, and
the only intersection with $L$ is at $(0,\varepsilon)\in R_2$, with
multiplicity $|\det((0,1),(-1,-1))|=|0\cdot(-1)-1\cdot(-1)|=1$.  Total: $1$.

\emph{The direction $v=(1,-1)$.}  The vertex is
$(a+\varepsilon,a-\varepsilon)$ and the third ray is
$\{(s+2\varepsilon,s):s\leq a-\varepsilon\}$.  Intersecting with $R_1$ gives
$s=0$ and the point $(2\varepsilon,0)$, which lies in $R_1$; intersecting
with $R_2$ gives $s=-2\varepsilon$, hence $x_2<0$; intersecting with $R_0$
gives $x_1-x_2=2\varepsilon\neq0$ for the third ray, and $s>0$ for the other
two.  Multiplicity $|\det((1,0),(-1,-1))|=1$.  Total: $1$.

In all three cases the total is $1$, and the intersection point tends to the
origin as $\varepsilon\to0$.  The stable intersection is therefore
$L\cdot L'=1\cdot[(0,0)]$, supported at the vertex of $L$, which indeed lies
in $|L|\cap|L'|=R_0$.  The count matches the transverse count of
\Cref{ex:ch10-two-lines-generic}, as it must, since $L'$ can be moved from
$(a,a)$ to $(3,1)$ through translations.

\emph{The excluded direction.}  For $v=(-1,-1)$ the translate
$L'+\varepsilon v=L+(a-\varepsilon)(1,1)$ is again a diagonal translate, and
the intersection remains one-dimensional for every $\varepsilon$.  Thus
$v=(-1,-1)$, and likewise $v=(1,1)$, are not generic.  The nongeneric set is
the line $\mathbf R(1,1)$, a proper rational subspace; this is the
finite union of proper subspaces appearing in the proof of
\Cref{thm:tropical-stable-intersection}.  In the language of that proof,
$\tau=\{0\}$ is the codimension-one cell of $L$ at which the displacement
must cross a wall, and $H=N_{\tau,\mathbf R}+N_{R_0'}=\mathbf R(1,1)$ is the
hyperplane.  Crossing it in the direction of the class of $(1,0)$ picks out
the facets $R_1$ with $\lambda=1$; crossing in the opposite direction picks
out $R_2$ with $\lambda=-1$, after writing the primitive normals of $L$ at
its vertex modulo $H$:
\[
 (1,0)\equiv w,\qquad (0,1)\equiv-w,\qquad (-1,-1)\equiv0
 \pmod{\mathbf R(1,1)},
\]
where $w$ generates $N/(N\cap H)$.  The balancing relation
$1\cdot1+1\cdot(-1)+1\cdot0=0$ is exactly the equality of the two counts.
\end{example}

\begin{figure}[ht]
\centering
\begin{tikzpicture}[x=1cm,y=1cm,>=Stealth]
  \coordinate (O) at (0,0);
  \coordinate (A) at (1.3,1.3);
  \coordinate (B) at (1.9,1.3);
  \coordinate (P) at (0.6,0);

  % L
  \draw[very thick] (O)--(3.0,0);
  \draw[very thick] (O)--(0,3.0);
  \draw[very thick] (O)--(-1.9,-1.9);
  \fill (O) circle (1.6pt);
  \node[below left,font=\small] at (-0.05,-0.05) {$0$};
  \node[font=\small] at (3.35,0) {$L$};

  % L' = L + (a,a)
  \draw[thick,gray,dash pattern=on 3pt off 2pt] (A)--(3.6,1.3);
  \draw[thick,gray,dash pattern=on 3pt off 2pt] (A)--(1.3,3.6);
  \draw[thick,gray,dash pattern=on 3pt off 2pt] (A)--(-1.9,-1.9);
  \fill[gray] (A) circle (1.6pt);
  \node[gray,font=\small] at (3.95,1.3) {$L'$};

  % displaced L' + eps v, v = (1,0)
  \draw[thin,dash pattern=on 1pt off 1.6pt] (B)--(3.9,1.3);
  \draw[thin,dash pattern=on 1pt off 1.6pt] (B)--(1.9,3.6);
  \draw[thin,dash pattern=on 1pt off 1.6pt] (B)--(-1.4,-2.0);
  \fill (B) circle (1.2pt);
  \node[font=\scriptsize] at (2.55,3.35) {$L'+\varepsilon v$};

  % intersection point
  \fill (P) circle (2.2pt);
  \draw[->,thin] (1.15,-0.95)--(0.68,-0.12);
  \node[font=\scriptsize,anchor=north] at (1.15,-1.0)
    {multiplicity $1$ at $(\varepsilon,0)$};

  % overlapping ray highlight
  \draw[line width=3pt,black!18] (O)--(-1.9,-1.9);
  \draw[very thick] (O)--(-1.9,-1.9);
  \node[font=\scriptsize,anchor=east] at (-1.95,-1.75) {$|L|\cap|L'|$};
\end{tikzpicture}
\caption{Two tropical lines in degenerate position.  The supports share the
whole ray $\mathbf R_{\geq0}(-1,-1)$, so the set-theoretic intersection has
the wrong dimension.  After an arbitrarily small generic displacement the
two curves meet in one transverse point of multiplicity one, which tends to
the vertex of $L$.  Displacing in the direction $(0,1)$ instead produces the
point $(0,\varepsilon)$ on the vertical ray, again with multiplicity one.}
\label{fig:ch10-stable-lines}
\end{figure}

\begin{example}[A plane and a line in $\mathbf R^3$]
\label{ex:ch10-plane-line-r3}
Let $H\subset\mathbf R^3$ be the standard tropical plane of
\Cref{ex:ch10-standard-plane}, with weight one on each of its six cones, and
let $L_3\subset\mathbf R^3$ be the standard tropical line of
\Cref{ex:ch10-standard-line}, with rays $\mathbf R_{\geq0}e_i$,
$i=0,1,2,3$, and weight one.  Since every ray of $L_3$ is a face of a cone of
$H$, we have $L_3\subseteq H$, so $H\cap L_3=L_3$ has dimension $1$, whereas
the expected dimension is $2+1-3=0$.  Again we displace.

It is convenient to use the description of $H$ as the corner locus of
$F(x)=\max\{0,-x_1,-x_2,-x_3\}$, proved in
\Cref{ex:ch10-standard-plane}: a point $x$ lies in $H$ if and only if the
maximum defining $F(x)$ is attained at least twice.

\emph{The direction $v=(3,2,1)$.}  The displaced line has vertex
$p=\varepsilon(3,2,1)$ and rays $p+\mathbf R_{\geq0}e_i$.  Along
$p+te_1=(3\varepsilon+t,2\varepsilon,\varepsilon)$ we get
$F=\max\{0,-3\varepsilon-t,-2\varepsilon,-\varepsilon\}=0$, attained only by
the constant term, so this ray misses $H$; the same computation applies to
$p+te_2$ and $p+te_3$.  Along
\[
 p+te_0=(3\varepsilon-t,\,2\varepsilon-t,\,\varepsilon-t),
\]
the four candidates are $0$, $t-3\varepsilon$, $t-2\varepsilon$,
$t-\varepsilon$, whose maximum is $\max\{0,t-\varepsilon\}$; it is attained
twice exactly when $t=\varepsilon$, giving the point
$q=(2\varepsilon,\varepsilon,0)$.  At $q$ the maximum is attained by $0$ and
$-x_3$, so $q$ lies in the relative interior of
$\operatorname{cone}(e_1,e_2)=\{x_3=0,\ x_1,x_2\geq0\}$.  The intersection
is transverse, with tangent lattices $\mathbf Ze_1+\mathbf Ze_2$ and
$\mathbf Ze_0=\mathbf Z(1,1,1)$, and the multiplicity is
\[
 1\cdot1\cdot
 \bigl|\det\bigl((1,0,0),(0,1,0),(1,1,1)\bigr)\bigr|
 =1 .
\]
Total: $1$.

\emph{The direction $v=(1,2,3)$.}  Now $p=\varepsilon(1,2,3)$, the three rays
$p+te_i$ with $i\geq1$ again miss $H$, and along $p+te_0$ the maximum is
$\max\{0,t-\varepsilon\}$, attained twice at $t=\varepsilon$, giving
$q'=(0,\varepsilon,2\varepsilon)$.  At $q'$ the maximum is attained by $0$
and $-x_1$, so $q'$ lies in $\operatorname{cone}(e_2,e_3)$, and the
multiplicity is
$|\det((0,1,0),(0,0,1),(1,1,1))|=1$.  Total: $1$, in agreement with the
previous direction, although the intersection point has moved to a different
cone of $H$.

\emph{An excluded direction.}  For $v=e_0=(-1,-1,-1)$ the vertex is
$-\varepsilon(1,1,1)$ and along the ray $p+te_0$ one has
$x_1=x_2=x_3=-\varepsilon-t$, so $F=\varepsilon+t$ is attained by three
terms at once; the whole ray lies in $H$ and the intersection stays
one-dimensional.  As in the planar case, the nongeneric directions form a
proper closed subset, here containing the line $\mathbf Re_0$.

As $\varepsilon\to0$ the intersection point tends to the origin in both
generic computations, so $H\cdot L_3=1\cdot[(0,0,0)]$.
\end{example}

The three computations exhibit the two phenomena that
\Cref{thm:tropical-stable-intersection} must explain: the count is the same
for all generic displacement directions, and the support of the answer is
contained in---but usually much smaller than---the set-theoretic
intersection.  \Cref{ex:ch10-unbalanced} showed that the count fails without balancing;
the support inclusion is formal, holding for any weighted complexes,
since \eqref{eq:tropical-stable-weight} puts weight only on cells
contained in $\sigma\cap\sigma'$.

\subsection{The general construction}
\label{sec:ch10-general-stable}

For a generic vector $v$, the cycles $X$ and $Y+\varepsilon v$ meet
transversely for every sufficiently small $\varepsilon>0$.  The
\emph{stable intersection} $X\cdot Y$ is supported on the
cells of dimension $k+l-r$ in a suitable refinement of $X\cap Y$; these
may lie in a higher-dimensional part of the original intersection.
When $k+l<r$, the answer is the zero cycle. If $k+l\geq r$ and $\tau$ is a
candidate cell of dimension $k+l-r$, set
\begin{equation}
\label{eq:tropical-stable-weight}
 \omega_{X\cdot Y}(\tau)
 =\sum_{\sigma,\sigma'}
   \omega_X(\sigma)\omega_Y(\sigma')
   [N:N_\sigma+N_{\sigma'}],
\end{equation}
where the sum is over pairs of facets for which
$\tau\subset\sigma\cap\sigma'$ and
$\sigma\cap(\sigma'+\varepsilon v)\neq\varnothing$ for all sufficiently
small $\varepsilon>0$.  Locally finite complexes are handled on relatively
compact sets; the resulting local cycles agree on overlaps.

Two points about this definition are easy to misread.  The condition
``$\sigma\cap(\sigma'+\varepsilon v)\neq\varnothing$ for all sufficiently
small $\varepsilon>0$'' selects the pairs of facets that continue to meet
after the displacement; a pair whose intersection is destroyed by an
arbitrarily small push contributes nothing, even if $\sigma\cap\sigma'$ is
nonempty.  And the sum is indexed by pairs of facets, not by points: several
pairs may contribute to the same cell $\tau$, and their contributions add.
The condition is local along $\RelInt(\tau)$. For $x\in\RelInt(\tau)$,
it is equivalent to
$v\in\operatorname{cone}_x(\sigma)-\operatorname{cone}_x(\sigma')$.
If $v=a-b$ for feasible directions $a,b$, then
$x+\varepsilon a\in\sigma\cap(\sigma'+\varepsilon v)$ for all small
$\varepsilon>0$. Conversely, subtracting the two directions from $x$ to a
displaced intersection point gives this condition. Thus contributing pairs
have displaced intersections approaching every such $x$. Cells with zero
total weight are omitted.
\Cref{ex:ch10-two-lines-degenerate} exhibits both features, since exactly
one of the three pairs through the origin survives each
displacement.

\begin{theorem}
\label{thm:tropical-stable-intersection}
The weight in \eqref{eq:tropical-stable-weight} is independent of the generic
vector and of all polyhedral refinements.  It defines a balanced tropical
cycle.  Stable intersection is bilinear, commutative, associative, and has
$N_{\mathbf R}$ with weight one as unit.  Moreover,
$|X\cdot Y|\subset |X|\cap|Y|$.
\end{theorem}
\begin{proof}
We give the fan-displacement proof, following
\cite[\S4.3]{MR18Tropical}.  Refining either complex merely partitions the
relative interiors of its facets.  At a generic displaced intersection,
each contributing old pair contains a unique contributing refined pair, and
subdivision does not change either tangent lattice.  Hence the summand and
the total weight in \eqref{eq:tropical-stable-weight} are unchanged.

We next prove independence of $v$.  The set of nongeneric translation
vectors is contained in a finite union of proper rational affine subspaces.
Within one chamber of the complement, the intersecting pairs of facets do
not change.  It remains to cross one wall.  Quotienting by the common tangent
space reduces the change to the following local configuration: $Y$ is a
linear cycle with tangent lattice $N'$, while $X$ is a fan whose facets
$\sigma_i$ meet along one codimension-one linear face with tangent lattice
$N_\tau$.  Put $H=N_{\tau,\mathbf R}+N'_{\mathbf R}$.  If $H$ has codimension
at least two, a generic translate misses the configuration on both sides and
there is nothing to prove.  Otherwise $H$ is a hyperplane.  Choose a
primitive generator $w$ of $N/(N\cap H)$ and write
\[
 u_{\sigma_i/\tau}=\lambda_iw\pmod H,
 \qquad \lambda_i\in\mathbf Z.
\]
Balancing gives
\[
                         \sum_i\omega_X(\sigma_i)\lambda_i=0.
\]
For $\lambda_i\neq0$ there is an index $q$, independent of $i$, such that
\[
 [N:N_{\sigma_i}+N']
 =|\lambda_i|q.
\]
Here $q=[N\cap H:N_\tau+N']$. Both sides measure successively this index
and the index $|\lambda_i|$ in the rank-one quotient $N/(N\cap H)$.
Facets with $\lambda_i=0$ are missed by a generic translate; their
non-full-rank lattice sum is not assigned a finite index. A translate crossing the wall in
the $w$-direction meets exactly the facets with $\lambda_i>0$; crossing in
the opposite direction meets those with $\lambda_i<0$.  Their total
contributions are
\[
 q\sum_{\lambda_i>0}\omega_X(\sigma_i)\lambda_i
 \quad\text{and}\quad
 -q\sum_{\lambda_i<0}\omega_X(\sigma_i)\lambda_i,
\]
which are equal by the displayed balancing equation.  Any two generic
vectors can be joined by a path crossing finitely many walls, proving
independence.

The construction is local.  More precisely, quotienting by the tangent
lattice of a cell $\tau$ preserves every index in
\eqref{eq:tropical-stable-weight} and gives
\[
 \operatorname{Star}_{X\cdot Y}(\tau)
 =\operatorname{Star}_X(\tau)\cdot\operatorname{Star}_Y(\tau).
\]
We first check balancing for a transverse product.  A codimension-one cell of
that product has the form $\rho=\tau\cap\sigma'$, where $\tau$ is a
codimension-one cell of $X$ and $\sigma'$ is a facet of $Y$, or the same
description with the factors exchanged.  Let $\sigma_i\supset\tau$ be the
facets of $X$, let $u_i$ be their primitive normal classes, and let $w_i$ be
the primitive normal classes of
$\sigma_i\cap\sigma'$ along $\rho$.  From the definitions of primitive
generators and multiplicativity of lattice indices,
\[
 [N:N_{\sigma_i}+N_{\sigma'}]\,w_i
 =[N:N_\tau+N_{\sigma'}]\,u_i
 \pmod {N_\tau}.
\]
Indeed, the coefficient relating $w_i$ to $u_i$ is
\[
 [N_{\sigma_i}:
   (N_{\sigma_i}\cap N_{\sigma'})+N_\tau],
\]
and multiplying it by $[N:N_{\sigma_i}+N_{\sigma'}]$ gives the common index
$[N:N_\tau+N_{\sigma'}]$.  Consequently the weighted sum of the $w_i$ is,
modulo $N_\tau$, the common factor
$\omega_Y(\sigma')[N:N_\tau+N_{\sigma'}]$ times
$\sum_i\omega_X(\sigma_i)u_i$, hence is zero modulo $N_\tau$.  The same sum
lies in $N_{\sigma'}$, because every $w_i$ does.  It therefore lies in
$N_\tau\cap N_{\sigma'}=N_\rho$, which is precisely the balancing condition
for the transverse product.

To check balancing of the stable product, quotient by a codimension-one cell
and use the star identity.  We are reduced to a one-dimensional fan obtained
as the recession fan of a generic transverse product.  Truncate that
one-cycle by a large polytope and sum the balancing equations at all its
vertices.  Contributions from bounded edges cancel in opposite pairs,
leaving the weighted sum of the recession rays equal to zero.  The weights
of those rays are exactly the stable weights in
\eqref{eq:tropical-stable-weight}.  Hence the stable product is balanced.

Bilinearity follows from the sum in
\eqref{eq:tropical-stable-weight}; commutativity follows by exchanging the
two factors.  The full affine space has tangent lattice $N$, so its
intersection index with any facet is one, proving the unit assertion.
The definition also gives the support inclusion.

It remains to prove associativity.  First choose simultaneous generic
translations so that all pairwise and triple intersections are transverse.
For tangent lattices $A,B,C\subset N$ of three participating facets, the
weight in $(X\cdot Y)\cdot Z$ contains
\[
 [N:A+B]\,[N:(A\cap B)+C].
\]
The elementary index identity
\[
 [N:(A\cap B)+C]
 =[N:B+C]\,[B:(A\cap B)+(B\cap C)]
\]
shows that the product equals
\[
 [N:A+B]\,[N:B+C]\,
 [B:(A\cap B)+(B\cap C)],
\]
which is unchanged when $A$ and $C$ are exchanged.  To verify the identity,
map $B$ into $N/((A\cap B)+C)$: its kernel is
$(A\cap B)+(B\cap C)$, and its cokernel is $N/(B+C)$; multiplicativity of
orders in this finite kernel--cokernel sequence gives the formula.
Thus the two parenthesizations have the same weight in the transverse case.
In general, compute both sides after two independent arbitrarily small
generic translations.  The wall-crossing independence already proved lets
both translations tend to zero, yielding
$(X\cdot Y)\cdot Z=X\cdot(Y\cdot Z)$.
\end{proof}

This is the convention and proof of
\cite[Propositions~4.3.6--4.3.7]{MR18Tropical}.  The CLD formulation is
\cite[\S6]{CLD25}. \Cref{chap:delta-forms} upgrades the product to
delta-forms and compares it with analytic divisor intersections.

It is worth naming the four uses of balancing in that proof, since they are
of different kinds.  Balancing was used once to prove independence of the
generic vector, in the wall-crossing step; that is where the hypothesis is
genuinely needed, and \Cref{ex:ch10-unbalanced} shows the conclusion fails
without it.  Balancing was used a second time to prove that the transverse
product is again balanced---a conclusion, not a hypothesis.  It was used a
third time, through the star identity, to reduce the balancing of the stable
product to a one-dimensional statement.  And it was used a fourth time in the
truncation argument, where summing the vertex equations of a bounded piece of
a one-cycle leaves only the recession contributions.  The reader who wants a
single mnemonic may take the following: balancing is exactly the condition
under which a weighted complex has no boundary, and a cycle without boundary
can be moved freely.

The following corollary is the form in which the theorem is applied in
\Cref{chap:delta-forms}.

\begin{corollary}
\label{cor:ch10-stable-star}
For tropical cycles $X,Y$ in $N_{\mathbf R}$ and a cell $\tau$ of a common
refinement contained in $|X|\cap|Y|$,
\[
 \operatorname{Star}_{X\cdot Y}(\tau)
 =\operatorname{Star}_X(\tau)\cdot\operatorname{Star}_Y(\tau)
\]
as tropical fan cycles in $N_{\mathbf R}/N_{\tau,\mathbf R}$, with the left
side zero if $\tau$ is not contained in the product support. In particular
the stable intersection is local: the germ of $X\cdot Y$ at a point depends
only on the germs of $X$ and $Y$ at that point.
\end{corollary}

\begin{proof}
The identity is proved as part of
\Cref{thm:tropical-stable-intersection}; we record why it gives locality.  A
germ at a point $x$ determines the star of the smallest cell containing $x$
in its relative interior, and conversely the star determines the germ up to
an integer-affine change that fixes $x$.  Since the displayed identity
expresses the star of the product in terms of the stars of the factors, two
pairs of cycles with the same germs at $x$ have products with the same germ
at $x$.
\end{proof}

\begin{example}[Two tropical lines]
\label{ex:ch10-two-lines-summary}
Take the tropical line in \Cref{fig:tropical-line-balancing} and a sufficiently
small generic translate.  A local pair of transverse rays has primitive
directions forming a lattice basis, so its determinant has absolute value
one.  The stable intersection is therefore one point of multiplicity one.
If one line is pulled back by the lattice map $(x,y)\mapsto(2x,y)$, the
relevant determinant can become two; this is exactly the index appearing in
the lattice test and in \Cref{prop:tropical-pushforward-balanced}.
\end{example}

\begin{remark}[Why ordinary smooth forms are not enough]
\label{rem:ch10-delta-forms-motivation}
A continuous piecewise affine function has well-defined slopes on cells, but
$\dprimeddprime$ of the function is supported on codimension-one faces and is not an
ordinary smooth superform. Delta-forms adjoin exactly these weighted
polyhedral currents. Harmonic tropicalizations, by contrast, enlarge the
allowed chart functions and improve the resulting sheaf cohomology. Both
extensions preserve the superform calculus developed here and are compared
systematically in \Cref{chap:delta-forms}.
\end{remark}

\section{Products, graphs, and Gysin maps}
\label{sec:ch10-gysin}
\index{Gysin map!tropical}\index{product!of tropical cycles}

Stable intersection was defined by a displacement.  There is a second,
equivalent, description, in which the displacement is replaced by a
pullback along the diagonal.  It has three advantages: it makes the
symmetry of the product manifest, it explains why lattice indices appear
multiplicatively, and it is the form that survives the passage to currents in
\Cref{chap:cld-currents,chap:delta-forms}, where displacement is unavailable.
The Chambert--Loir--Ducros construction of tropical intersection is of this
second kind, carried out at the level of positive closed currents; the
comparison is discussed at the end of the section.

\subsection{Products of cycles}

\begin{proposition}
\label{prop:ch10-product-cycle}
Let $(\mathscr C,m)$ and $(\mathscr C',m')$ be balanced weighted rational
polyhedral complexes in $N_{\mathbf R}$ and $N'_{\mathbf R}$, of pure
dimensions $n$ and $n'$.  Then the cells $\sigma\times\sigma'$ with the
weights $m_\sigma m'_{\sigma'}$ form a balanced weighted rational polyhedral
complex $\mathscr C\times\mathscr C'$ in $N_{\mathbf R}\oplus
N'_{\mathbf R}$, of pure dimension $n+n'$, with tangent lattices
$N_{\sigma}\oplus N'_{\sigma'}$.  The construction descends to tropical
cycles and is bilinear; we write $X\boxtimes Y$ for the resulting cycle.
\end{proposition}

\begin{proof}
A product of rational polyhedra is a rational polyhedron, its faces are the
products of faces, and the tangent space of a product is the product of the
tangent spaces; intersecting with the lattice $N\oplus N'$ gives
$N_\sigma\oplus N'_{\sigma'}$, which is saturated because each factor is.
Face-closedness and the intersection axiom follow factorwise, and the
dimension is additive.

For balancing, let $\rho$ be a codimension-one cell of the product.  Since
$\dim(\sigma\times\sigma')=\dim\sigma+\dim\sigma'$ and cells of the product
are products of cells, $\rho=\tau\times\sigma'$ with $\tau$ of codimension
one in $\mathscr C$ and $\sigma'$ a facet of $\mathscr C'$, or
$\rho=\sigma\times\tau'$ symmetrically.  Treat the first case.  The facets of
the product containing $\rho$ are the $\sigma\times\sigma'$ with
$\sigma\supset\tau$, their weights are $m_\sigma m'_{\sigma'}$, and the
primitive normal of $\sigma\times\sigma'$ along $\rho$ is the image of
$u_{\sigma/\tau}$ under
\[
 N_\sigma/N_\tau
 \hookrightarrow
 (N\oplus N')/(N_\tau\oplus N'_{\sigma'}),
\]
which is injective because $N_\sigma\cap(N_\tau\oplus N'_{\sigma'})$, computed
inside $N\oplus 0$, equals $N_\tau$.  Hence the balancing sum at $\rho$ is
$m'_{\sigma'}$ times the image of $\sum_{\sigma\supset\tau}m_\sigma
u_{\sigma/\tau}=0$, hence zero.  The second case is symmetric.  Bilinearity
and compatibility with refinement are immediate, since a common refinement of
the factors induces one of the product.
\end{proof}

\begin{proposition}
\label{prop:ch10-graph-cycle}
Let $F:N_{\mathbf R}\to N'_{\mathbf R}$ be integer affine with linear part
$A$, and let
\[
 \Gamma_F=\{(x,F(x))\mid x\in N_{\mathbf R}\}
 \subset N_{\mathbf R}\oplus N'_{\mathbf R}.
\]
Then $\Gamma_F$ is a translate of the real span of the saturated sublattice
$N_{\Gamma}=\{(a,A(a)):a\in N\}$ of $N\oplus N'$, and, with weight one, it is
a tropical cycle of dimension $\operatorname{rank}N$.  The projection
$p:N_{\mathbf R}\oplus N'_{\mathbf R}\to N_{\mathbf R}$ restricts to an
integer-affine isomorphism $\Gamma_F\to N_{\mathbf R}$ carrying $N_\Gamma$
isomorphically onto $N$.
\end{proposition}

\begin{proof}
$\Gamma_F$ is the translate by $(0,F(0))$ of the graph of $A$, which is the
kernel of the surjection $N_{\mathbf R}\oplus N'_{\mathbf R}\to
N'_{\mathbf R}$, $(a,b)\mapsto b-A(a)$.  A kernel of a lattice homomorphism
is saturated: if $k(a,b)\in N_\Gamma$ then $kb=kA(a)$, so $b=A(a)$.  A
translate of a saturated sublattice, with weight one, is balanced by
\Cref{cor:ch10-balancing-consequences}\,\cref{itm:ch10-bal-affine}.  The
projection $p$ restricted to $\Gamma_F$ has the integer-affine inverse
$x\mapsto(x,F(x))$, and it carries $(a,A(a))$ to $a$, hence maps $N_\Gamma$
isomorphically onto $N$.
\end{proof}

\subsection{The Gysin pullback}

\begin{definition}
\label{def:ch10-gysin}
Let $F:N_{\mathbf R}\to N'_{\mathbf R}$ be integer affine, with graph
inclusion $j:N_{\mathbf R}\to N_{\mathbf R}\oplus N'_{\mathbf R}$,
$j(x)=(x,F(x))$, and first projection $p$.  For a tropical cycle $S$ on
$N_{\mathbf R}\oplus N'_{\mathbf R}$ put
\[
 j^!S=p_*\bigl(S\cdot[\Gamma_F]\bigr),
\]
where $\cdot$ is the stable intersection of
\Cref{thm:tropical-stable-intersection} and $p_*$ is the pushforward of
\Cref{sec:ch10-pushforward}.
\end{definition}

\begin{lemma}
\label{lem:ch10-gysin-defined}
The right-hand side of \Cref{def:ch10-gysin} is defined, and $p_*$ introduces
no lattice index: if $\rho$ is a facet of $S\cdot[\Gamma_F]$ then
$[N_{p(\rho)}:p(N_\rho)]=1$ and $p$ is injective on $\rho$.  Consequently
$j^!S$ is zero when $\dim S<\operatorname{rank}N'$, and otherwise is a
tropical cycle of dimension $\dim S-\operatorname{rank}N'$.
\end{lemma}

\begin{proof}
By \Cref{thm:tropical-stable-intersection}, $S\cdot[\Gamma_F]$ is zero when
the expected dimension is negative.  Otherwise it is a tropical cycle
supported in $|S|\cap\Gamma_F\subseteq\Gamma_F$, of dimension
$\dim S+\operatorname{rank}N-(\operatorname{rank}N+\operatorname{rank}N')
=\dim S-\operatorname{rank}N'$.  Every facet $\rho$ of it lies in
$\Gamma_F$, so $N_\rho\subseteq N_\Gamma$ and $p|_\rho$ is injective, being
the restriction of the isomorphism $p|_{\Gamma_F}$ of
\Cref{prop:ch10-graph-cycle}.  That isomorphism carries $N_\Gamma$ onto $N$,
hence carries the saturated subgroup $N_\rho\subseteq N_\Gamma$ onto the
saturated subgroup $p(N_\rho)=N_{p(\rho)}\subseteq N$, so the index is one.
The restriction of $p$ to the closed support inside $\Gamma_F$ is proper,
since $p|_{\Gamma_F}$ is a homeomorphism onto $N_{\mathbf R}$.
\Cref{prop:tropical-pushforward-balanced} now applies.
\end{proof}

\begin{theorem}
\label{thm:ch10-diagonal-gysin}
Let $X$ and $Y$ be tropical cycles in $N_{\mathbf R}$, let
$\delta:N_{\mathbf R}\to N_{\mathbf R}\oplus N_{\mathbf R}$ be the diagonal
$\delta(x)=(x,x)$, and let $j=\delta$ be the graph inclusion of
$F=\mathrm{id}$.  Then
\[
 X\cdot Y=j^!\,(X\boxtimes Y).
\]
More generally, for an integer-affine $F:N_{\mathbf R}\to N'_{\mathbf R}$ and
cycles $X$ on $N_{\mathbf R}$, $Z$ on $N'_{\mathbf R}$, the cycle
$j^!(X\boxtimes Z)$ is supported in $|X|\cap F^{-1}(|Z|)$ and equals
$X\cdot F^*Z$ whenever $F$ is an isomorphism, where
$F^*Z=(F^{-1})_*Z$.
\end{theorem}

\begin{proof}
\emph{Step 1: displacing the diagonal.}
By the commutativity in \Cref{thm:tropical-stable-intersection} we may
displace either factor of $(X\boxtimes Y)\cdot[\Delta]$; we displace the
diagonal.  Fix a generic $v\in N_{\mathbf R}$ and replace $\Delta$ by
$\Delta+\varepsilon(0,-v)$.  Then
\[
 (X\boxtimes Y)\cap\bigl(\Delta+\varepsilon(0,-v)\bigr)
 =\{(x,x-\varepsilon v)\mid x\in|X|,\ x-\varepsilon v\in|Y|\},
\]
which is carried by $p$ onto $|X|\cap(|Y|+\varepsilon v)$.  Thus the
displaced set-theoretic intersections on the two sides correspond, and a
vector $v$ is generic for one problem exactly when it is generic for the
other, since the two families of walls correspond under the same
isomorphism.

\emph{Step 2: matching the multiplicities.}
Let $\sigma$ be a facet of $X$ and $\sigma'$ a facet of $Y$ whose displaced
intersection is nonempty and transverse.  The corresponding facet of
$X\boxtimes Y$ is $\sigma\times\sigma'$ with weight
$\omega_X(\sigma)\omega_Y(\sigma')$ and tangent lattice
$N_\sigma\oplus N_{\sigma'}$, by \Cref{prop:ch10-product-cycle}.  Its
transverse multiplicity against $\Delta$, whose weight is one and whose
tangent lattice is $N_\Delta=\{(a,a)\}$, is
\[
 \omega_X(\sigma)\omega_Y(\sigma')
 \bigl[N\oplus N:(N_\sigma\oplus N_{\sigma'})+N_\Delta\bigr]
 =\omega_X(\sigma)\omega_Y(\sigma')\,[N:N_\sigma+N_{\sigma'}],
\]
by \Cref{lem:ch10-index}\,\cref{itm:ch10-idx-diagonal}.  This is exactly the
summand of \eqref{eq:tropical-stable-weight}.

\emph{Step 3: the projection.}
By \Cref{lem:ch10-gysin-defined}, $p_*$ neither changes weights nor merges
distinct facets, since $p|_\Delta$ is an isomorphism.  Hence, letting
$\varepsilon\to0$, the two constructions produce the same weight on every
facet of the limit, and by
\Cref{thm:tropical-stable-intersection} both limits exist and are independent
of $v$.  Therefore $X\cdot Y=j^!(X\boxtimes Y)$.

\emph{Step 4: the general graph.}
For arbitrary $F$, the set $|X\boxtimes Z|\cap\Gamma_F$ consists of the
points $(x,F(x))$ with $x\in|X|$ and $F(x)\in|Z|$, whose image under $p$ is
$|X|\cap F^{-1}(|Z|)$; the support statement follows from
$|S\cdot T|\subseteq|S|\cap|T|$.  If $F$ is an isomorphism, meaning that its
linear part carries $N$ isomorphically onto $N'$, then the integer-affine
isomorphism
$(\mathrm{id}\oplus F^{-1}):N_{\mathbf R}\oplus N'_{\mathbf R}\to
N_{\mathbf R}\oplus N_{\mathbf R}$ carries $\Gamma_F$ to $\Delta$ and
$X\boxtimes Z$ to $X\boxtimes F^*Z$, both integer-linearly and preserving all
weights and lattices by
\Cref{prop:tropical-pushforward-balanced}; Steps 1--3 then give
$j^!(X\boxtimes Z)=X\cdot F^*Z$.
\end{proof}

\begin{example}[The two tropical lines, recomputed on the diagonal]
\label{ex:ch10-lines-via-diagonal}
Take $X=L$ and $Y=L'=L+(3,1)$ as in
\Cref{ex:ch10-two-lines-generic}.  The product $X\boxtimes Y$ is a
two-dimensional cycle in $\mathbf R^4$ with nine facets
$R_i\times R_j'$, all of weight one.  The diagonal
$\Delta\subset\mathbf R^2\oplus\mathbf R^2$ is two-dimensional, so the
expected dimension of $(X\boxtimes Y)\cdot\Delta$ is $2+2-4=0$.  The unique
facet meeting $\Delta$ is $R_1\times R_0'$, at the point
$((2,0),(2,0))$, and the multiplicity is
\[
 \bigl[\mathbf Z^4:(\mathbf Z(1,0)\oplus\mathbf Z(-1,-1))+N_\Delta\bigr]
 =\bigl[\mathbf Z^2:\mathbf Z(1,0)+\mathbf Z(-1,-1)\bigr]=1,
\]
by \Cref{lem:ch10-index}\,\cref{itm:ch10-idx-diagonal}.  Projecting back gives
the point $(2,0)$ with multiplicity one, in agreement with the direct
computation.  The gain is conceptual rather than computational: the
symmetric roles of $X$ and $Y$ are now built into the construction, and the
only lattice index that occurs is the diagonal index of
\Cref{lem:ch10-index}.
\end{example}

\begin{proposition}
\label{prop:ch10-gysin-projection}
In the setting of \Cref{def:ch10-gysin}, let $S$ be a tropical cycle on
$N_{\mathbf R}\oplus N'_{\mathbf R}$ and let $\beta$ be a smooth superform on
$N_{\mathbf R}$ of bidegree $(d,d)$, where
$d=\dim S-\operatorname{rank}N'\geq0$, with $j^!S$ of compact support.
Then
\[
 \int_{j^!S}\beta
 =\int_{S\cdot[\Gamma_F]}p^*\beta .
\]
\end{proposition}

\begin{proof}
This is the change-of-variables statement of
\Cref{chap:tropical-integration} applied to the integer-affine map $p$,
which by \Cref{lem:ch10-gysin-defined} is injective on every facet of
$S\cdot[\Gamma_F]$ with lattice index one.  Properness holds because the
support of $S\cdot[\Gamma_F]$ lies in $\Gamma_F$, on which $p$ is a
homeomorphism, and $j^!S$ was assumed compactly supported.
\end{proof}

\begin{warning}{Two Gysin maps}
The map $j^!$ above is defined on \emph{cycles} by displacement.  The
Chambert--Loir--Ducros Gysin map is defined on positive closed
\emph{currents}, by intersecting with the graph through iterated
Bedford--Taylor products of the parlinear functions
$u_i(x,y)=\max\{F_i(x),y_i\}$ and then pushing forward; see
\cite[\S9.5]{CLD25}.  The two constructions agree on cycles, and the
comparison is carried out in \Cref{chap:delta-forms} after currents are
available.  Until then, no statement about currents should be inferred from
the cycle-level construction, and in particular the harmonicity hypotheses
that CLD impose on $F$ are automatic for the integer-affine coordinate
functions used here, but must be checked for more general maps.
\end{warning}

\section{Dictionary}
\label{sec:ch10-dictionary}

The following table records what has been established, in the form in which
it is used later.  The left column is the object introduced here; the middle
column is the property that makes it usable; the right column is the first
place of use.

\begin{center}
\small
\begin{tabular}{p{.29\textwidth}|p{.33\textwidth}|p{.28\textwidth}}
\textbf{Object of this chapter}&\textbf{Established property}
  &\textbf{First use}\\ \hline
tangent lattice $N_\sigma$
  &saturated, of rank $\dim\sigma$
   (\cref{lem:ch10-tangent-lattice})
  &\cref{prop:integral-subdivision-independent}\\
primitive normal $u_{\sigma/\tau}$
  &canonical in $N_\sigma/N_\tau$, insensitive to lifts
   (\cref{ex:ch10-primitive-quotient})
  &\eqref{eq:superform-boundary-traces}\\
weight $m_\sigma$
  &unchanged by subdivision
   (\cref{lem:ch10-subdivision-lattices})
  &\cref{prop:integral-subdivision-independent}\\
balancing
  &vanishing of $\partial\mu$
   (\cref{thm:ch10-balancing-boundary})
  &\cref{thm:balanced-stokes}\\
balancing, second form
  &equivalent to Stokes on a fan
   (\cref{thm:ch10-balancing-stokes-fan})
  &\cref{thm:cld-global-stokes}\\
star of a cell
  &balanced, and computes local data
   (\cref{cor:ch10-balancing-consequences})
  &\cref{cor:ch10-stable-star}\\
$A^{p,q}(\Omega)$
  &sheaf, invariant under the stated chart changes
   (\cref{thm:ch10-forms-presentation-independent})
  &\cref{chap:analytic-tropicalization}\\
tropical space
  &transitions transport forms, lattices, calibration
   (\cref{thm:ch10-transition-transport})
  &\cref{thm:cld-global-stokes}\\
calibration $\mu_\sigma$
  &oriented multiplicity in the top exterior lattice
   (\cref{prop:ch10-calibration-basics})
  &\cref{prop:appG-calibration-tests}\\
pushforward index
  &$[N'_{\sigma'}:F(N_\sigma)]$, multiplicative
   (\cref{prop:tropical-pushforward-balanced})
  &\cref{thm:appG-analytic-calibration}\\
stable intersection
  &independent of displacement; local
   (\cref{thm:tropical-stable-intersection},
    \cref{cor:ch10-stable-star})
  &\cref{thm:delta-preform-algebra}\\
diagonal description
  &$X\cdot Y=j^!(X\boxtimes Y)$
   (\cref{thm:ch10-diagonal-gysin})
  &\cref{chap:delta-forms}
\end{tabular}
\end{center}

\subsection{From the general theory back to this chapter}

The reverse dictionary lists, for each later abstraction, the smallest
instance computed here.

\begin{center}
\small
\begin{tabular}{p{.30\textwidth}|p{.32\textwidth}|p{.28\textwidth}}
\textbf{Later object}&\textbf{Smallest instance}&\textbf{Computed in}\\ \hline
calibrated skeleton of an analytic space
  &tropical line $L_2$, weights $1$
  &\cref{ex:ch10-standard-line}\\
balanced tropicalization of a hypersurface
  &standard plane $H\subset\mathbf R^3$
  &\cref{ex:ch10-standard-plane}\\
balanced with realizability left open
  &five-ray fan $W\subset\mathbf R^3$
  &\cref{ex:ch10-codim-two-fan}\\
failure of Stokes without harmony
  &$\mathbf R_{\geq0}e_1\cup\mathbf R_{\geq0}e_2$
  &\cref{ex:ch10-unbalanced}\\
matroidal fan
  &tripod $B(U_{2,3})$
  &\cref{ex:ch10-bergman-u23}\\
abstract chart, no effective balanced embedding
  &circle $\mathbf R/\ell\mathbf Z$
  &\cref{ex:ch10-circle}\\
tropical multiplicity
  &$|\det((1,0),(1,2))|=2$
  &\cref{ex:ch10-lattice-index-test}\\
stable intersection of currents
  &two tropical lines, one point
  &\cref{ex:ch10-two-lines-generic},
   \cref{ex:ch10-two-lines-degenerate}\\
intersection in codimension two
  &$H\cdot L_3=[0]$ in $\mathbf R^3$
  &\cref{ex:ch10-plane-line-r3}\\
degree of a finite morphism
  &$[\mathbf Z:2\mathbf Z]=2$
  &\cref{ex:ch10-degree-two}
\end{tabular}
\end{center}

\subsection{Normalization checks}

The following computations have been carried out explicitly and should be
reproducible by the reader; together they pin down every convention of the
chapter.

\begin{enumerate}
\item $\sum_{i=0}^re_i=0$ is the balancing equation of the standard tropical
      line, and of the standard tropical plane on each of its rays
      (\cref{ex:ch10-standard-line,ex:ch10-standard-plane}).
\item $2(1,0)+3(0,1)+(-2,-3)=0$: balancing does not require equal weights
      (\cref{ex:ch10-weights-234}).
\item $\int_X\dprime\eta=-\langle a(0),\sum_km_ku_k\rangle$, which equals
      $-1$ for the unbalanced pair of rays
      (\cref{thm:ch10-balancing-stokes-fan},
       \cref{ex:ch10-unbalanced}).
\item $|\det((1,0),(1,2))|=2$ and $|\det((1,0),(-1,-1))|=1$: the two
      transverse multiplicities used throughout
      (\cref{ex:ch10-lattice-index-test,ex:ch10-two-lines-generic}).
\item $\bigl[\mathbf Z^4:(A\oplus B)+N_\Delta\bigr]=[\mathbf Z^2:A+B]$: the
      diagonal index
      (\cref{lem:ch10-index}\,\cref{itm:ch10-idx-diagonal}).
\item Two tropical lines meet with total multiplicity $1$, in general and in
      degenerate position, for every generic displacement
      (\cref{ex:ch10-two-lines-generic,ex:ch10-two-lines-degenerate}).
\item A tropical plane and a tropical line in $\mathbf R^3$ meet with total
      multiplicity $1$ (\cref{ex:ch10-plane-line-r3}).
\item Reversing the orientation of a cell, with the weight held fixed, flips
      $e_\sigma$, $\mu_\sigma$ and $[\sigma:\tau]$, and leaves the integral
      of an $(n,n)$-form and the calibrated balancing equation unchanged;
      flipping one weight in the displayed tripod breaks balancing
      (\cref{ex:ch10-sign-test-segment,ex:ch10-sign-test-line}).
\end{enumerate}

\section*{Source notes}

Our terms \emph{integer affine}, \emph{rational polyhedron}, \emph{facet},
\emph{tropical cycle}, and \emph{stable intersection}, as well as the
lattice-index formula for pushforward, follow
\cite[\S\S1.3, 2.1, 4.2--4.3 and 4.6]{MR18Tropical}.  The abstract
max-plus coordinate $x$ is related to the analytic valuation coordinate by
$x=-u$.  Calibrated tropical and parlinear spaces follow
\cite[\S\S3--6]{CLD25}. The later delta-form interface follows
Gubler--Künnemann, while harmonic-chart refinements follow
Gubler--Jell--Rabinoff; their exact bibliography keys are recorded in
\Cref{chap:delta-forms}.

\subsection*{Provenance of the individual statements}

The convex-geometric facts of \Cref{sec:ch10-polyhedra} are classical; the
two standard references for the polyhedral part are \cite{Roc70} and
\cite{Sch14}, and the lattice-theoretic part is elementary.  Nothing there is
special to tropical geometry, and the only nonstandard convention is that a
tangent lattice is always saturated
(\Cref{lem:ch10-tangent-lattice}); this is what makes primitive normal
vectors well defined and is the point at which a careless clearing of
denominators produces wrong multiplicities.

Balancing goes back, in the guise of the vanishing of a residue, to the
earliest tropical literature and appears in essentially the present form in
\cite[\S4.2]{MR18Tropical}, from which \Cref{def:ch10-weighted-balanced} is
adapted to our notation. The reformulation of balancing as the vanishing of a
boundary (\Cref{thm:ch10-balancing-boundary}) is the elementary shadow of the
chain-level description used in tropical homology, developed by Itenberg,
Katzarkov, Mikhalkin, and Zharkov; we have proved only the piece required
here and have deliberately not built the surrounding homology theory.  The
equivalence between balancing and Stokes' formula on a fan
(\Cref{thm:ch10-balancing-stokes-fan}) is a calculation, but the point of
view---that balancing is exactly the harmony of a calibration---is that of
Chambert--Loir and Ducros, whose condition d'harmonie is
\cite[Theorem~6.6.8 and \S15.6]{CLD25}; the analytic form of the equivalence
is \Cref{thm:appG-analytic-calibration}.

The Bergman fan of a matroid (\Cref{def:ch10-bergman-fan}) was introduced by
Sturmfels, and the description by chains of flats used here, together with
the general balancing theorem, is due to Ardila and Klivans.  We prove only
the rank-two case (\Cref{prop:ch10-bergman-rank-two}), which suffices for the
anchor example $U_{2,3}$ and is entirely self-contained; the general theorem
is not used anywhere in this book, and no bibliography key is assigned to it
here.  The reader who wants the systematic treatment, including the
identification of $B(M)$ with the tropicalization of a realization and the
relation to the Chow ring of a matroid, should consult the tropical-geometry
literature; \cite{MR18Tropical} is again the closest systematic reference in
our bibliography, and its \S4.2 contains the fan examples we have adapted.

The stable intersection product and the fan-displacement proof are
\cite[\S4.3]{MR18Tropical}; the same product occurs in the toric literature
under the name of the fan displacement rule.  The identification of stable
intersection with a Gysin pullback along the diagonal
(\Cref{thm:ch10-diagonal-gysin}) is the cycle-level shadow of the
current-level construction of \cite[\S\S9.5--9.6]{CLD25}, which builds the
product from iterated Bedford--Taylor products of the parlinear functions
$\max\{F_i(x),y_i\}$; that construction requires harmonicity hypotheses that
are automatic in our integer-affine setting.  Chambert--Loir and Ducros note
that the same product was obtained independently by other authors, and that
their own approach follows the tropical-intersection literature; we have kept
the elementary route because the book needs the product only on cycles until
\Cref{chap:delta-forms}.  The delta-form upgrade, in which the product is
extended to forms with polyhedral coefficients, is
\cite{GK17,GK19}; the harmonic-tropicalization enlargement of the chart
category is \cite{GJR26}; the toric global theory that uses these
multiplicities systematically is \cite{BGJKK25}.

Two omissions are deliberate.  We have not developed the theory of tropical
cycles with rational or real weights, which is what analytic calibrations
produce in the general Chambert--Loir--Ducros setting
(see the discussion after \eqref{eq:appG-analytic-multiplicity}); integrality
is an extra hypothesis, and every statement above that mentions a lattice
index has been formulated so that it survives the passage to real weights.
And we have not discussed realizability: the question of which balanced
cycles are tropicalizations is, as \Cref{ex:ch10-codim-two-fan} was designed
to make visible, not settled by anything in this chapter, and it plays no
role in the sequel, where the tropicalizations are given and the balancing is
a theorem about them rather than a characterization of them.
